\documentclass[a4paper,oneside]{report}
\usepackage[fontsize=13pt]{scrextend}      % Cỡ chữ 13pt (Normal của mẫu)
\usepackage[utf8]{vietnam}
\usepackage{mathptmx}                      % Times (tương đương Times New Roman)
\usepackage{amsmath}
\usepackage{graphicx}
\usepackage[table]{xcolor}
\usepackage{geometry}
\usepackage{setspace}
\usepackage{titlesec}
\usepackage[titles]{tocloft}
\usepackage{caption}
\usepackage{fancyhdr}
\usepackage{array}
\usepackage{tabularx}
\usepackage{longtable}                     % Bảng dài nhiều trang (danh mục ký hiệu, ma trận truy vết)
\usepackage{booktabs}                      % Kẻ bảng \toprule, \midrule, \bottomrule
\usepackage{multirow}
\usepackage{float}                         % Cho phép đặt hình/bảng cố định [H]
\usepackage{enumitem}                      % Chỉnh khoảng cách danh sách gạch đầu dòng
\usepackage{pdfpages}                      % Chèn bản scan Nhận xét GVHD, Phiếu báo cáo tiến độ
\usepackage{listings}                      % Mã nguồn ở Phụ lục A
\usepackage{makeidx}
\usepackage{tikz}
\usepackage{adjustbox}
\usetikzlibrary{shapes.geometric,arrows.meta,positioning,fit,backgrounds,calc}
\makeindex

% ================= LỀ TRANG =================
% Theo mẫu ĐATN: A4, trên 2cm, dưới 2cm, trái 3.5cm, phải 2.5cm
\geometry{
    a4paper,
    top=2cm,
    bottom=2cm,
    left=3.5cm,
    right=2.5cm,
    footskip=0.9cm
}

% ================= ĐOẠN VĂN =================
% Mẫu Word: Times 13pt, giãn dòng "Multiple 1.1", cách trước 3pt, căn đều hai bên,
% không thụt đầu dòng. Dòng cơ sở của LaTeX đã rộng hơn Word (~1.2 lần cỡ chữ),
% nên \setstretch{1.05} cho khoảng cách dòng tương đương "Multiple 1.1" của Word.
\setstretch{1.05}
\setlength{\parskip}{3pt}
\setlength{\parindent}{0pt}
\setlength{\emergencystretch}{2em}   % tránh chữ tràn lề

% Danh sách gạch đầu dòng: thụt 1.27cm như "List Paragraph" của mẫu, không giãn dòng thừa
\setlist{leftmargin=1.27cm, topsep=0pt, itemsep=0pt, parsep=3pt}

% Bảng: giãn hàng cho dễ đọc
\renewcommand{\arraystretch}{1.2}

% ================= SỐ TRANG =================
% Số trang ở chân trang, căn phải
\pagestyle{fancy}
\fancyhf{}
\fancyfoot[R]{\thepage}
\renewcommand{\headrulewidth}{0pt}
\fancypagestyle{plain}{%
    \fancyhf{}%
    \fancyfoot[R]{\thepage}%
    \renewcommand{\headrulewidth}{0pt}%
}

% ================= TIÊU ĐỀ CÁC CẤP =================
% Tối đa 3 cấp đánh số: 1.1 / 1.1.1 / 1.1.1.1 (cấp 1.1.1.1 in nghiêng).
% Cần thêm cấp nữa thì dùng a), b), c) hoặc gạch đầu dòng.
\setcounter{secnumdepth}{3}
\setcounter{tocdepth}{2}

% Chương: "CHƯƠNG 1. TÊN CHƯƠNG" – in hoa, đậm, căn giữa, 13pt
\titleformat{\chapter}[block]
    {\centering\normalsize\bfseries}
    {CHƯƠNG \thechapter.}{0.5em}{\MakeUppercase}
\titleformat{name=\chapter,numberless}[block]
    {\centering\normalsize\bfseries}
    {}{0pt}{\MakeUppercase}
\titlespacing*{\chapter}{0pt}{6pt}{12pt}

% Cấp 1 (1.1) – đậm
\titleformat{\section}[hang]
    {\normalsize\bfseries}{\thesection}{1em}{}
\titlespacing*{\section}{0pt}{6pt}{0pt}

% Cấp 2 (1.1.1) – đậm
\titleformat{\subsection}[hang]
    {\normalsize\bfseries}{\thesubsection}{1em}{}
\titlespacing*{\subsection}{0pt}{6pt}{0pt}

% Cấp 3 (1.1.1.1.) – nghiêng
\titleformat{\subsubsection}[hang]
    {\normalsize\itshape}{\thesubsubsection.}{1em}{}
\titlespacing*{\subsubsection}{0pt}{2pt}{0pt}

% ================= MỤC LỤC & DANH MỤC =================
\renewcommand{\cftchapfont}{\bfseries}
\renewcommand{\cftchappagefont}{\bfseries}
\renewcommand{\cftchapleader}{\cftdotfill{\cftdotsep}}
\renewcommand{\cftchappresnum}{CHƯƠNG }
\renewcommand{\cftchapaftersnum}{.}
\setlength{\cftchapnumwidth}{2.7cm}
\setlength{\cftbeforechapskip}{6pt}

\setlength{\cftsecindent}{0pt}
\setlength{\cftsecnumwidth}{1cm}
\setlength{\cftbeforesecskip}{4pt}

\setlength{\cftsubsecindent}{1cm}
\setlength{\cftsubsecnumwidth}{1.3cm}
\setlength{\cftbeforesubsecskip}{4pt}

% Danh mục hình vẽ: "Hình 2.1 Tên hình"
\renewcommand{\cftfigpresnum}{Hình }
\setlength{\cftfigindent}{0pt}
\setlength{\cftfignumwidth}{2cm}
\setlength{\cftbeforefigskip}{4pt}

% Danh mục bảng biểu: "Bảng 1.1 Tên bảng"
\renewcommand{\cfttabpresnum}{Bảng }
\setlength{\cfttabindent}{0pt}
\setlength{\cfttabnumwidth}{2cm}
\setlength{\cftbeforetabskip}{4pt}

% ================= CHÚ THÍCH HÌNH / BẢNG =================
% Theo style "Caption" của mẫu: nghiêng, 12pt, màu xám 404040, căn giữa, giãn dòng đơn.
% Hình: chú thích dưới; bảng: chú thích trên.
% Mẫu khuyến nghị hình/bảng rộng không quá 75% bề rộng phần chữ (width=0.75\linewidth).
\DeclareCaptionFont{capsize}{\fontsize{12}{14}\selectfont}
\DeclareCaptionFont{capcolor}{\color[HTML]{404040}}
\captionsetup{
    font={capsize,it,capcolor,singlespacing},
    labelfont={capsize,it,capcolor},
    labelsep=space,
    justification=centering,
    singlelinecheck=false,
    skip=6pt
}
\captionsetup[figure]{position=bottom}
\captionsetup[table]{position=top}

% ================= PHƯƠNG TRÌNH: "PT 1.1" =================
\numberwithin{equation}{chapter}
\makeatletter
\renewcommand{\tagform@}[1]{\maketag@@@{\ignorespaces PT #1\unskip\@@italiccorr}}
\makeatother

% ================= MÃ NGUỒN (PHỤ LỤC A) =================
% Định nghĩa ngôn ngữ Solidity cho listings
\lstdefinelanguage{Solidity}{
    morekeywords={pragma,solidity,contract,interface,library,is,function,modifier,event,emit,
        struct,enum,mapping,public,private,internal,external,view,pure,payable,returns,return,
        memory,storage,calldata,constructor,override,virtual,require,revert,if,else,for,while,
        new,delete,import,using,address,bool,string,bytes,bytes32,uint,uint8,uint256,int,int256},
    sensitive=true,
    morecomment=[l]{//},
    morecomment=[s]{/*}{*/},
    morestring=[b]",
}
\lstset{
    basicstyle=\ttfamily\footnotesize\singlespacing,
    keywordstyle=\bfseries,
    commentstyle=\itshape\color{gray},
    numbers=left,
    numberstyle=\tiny,
    frame=single,
    breaklines=true,
    columns=fullflexible,
    keepspaces=true,
    showstringspaces=false,
    tabsize=4
}
\renewcommand{\lstlistingname}{Đoạn mã}

% ================= TÊN CÁC MỤC TIẾNG VIỆT =================
\AtBeginDocument{%
    \renewcommand{\contentsname}{Mục lục}%
    \renewcommand{\listfigurename}{Danh mục hình vẽ}%
    \renewcommand{\listtablename}{Danh mục bảng biểu}%
    \renewcommand{\bibname}{Tài liệu tham khảo}%
    \renewcommand{\figurename}{Hình}%
    \renewcommand{\tablename}{Bảng}%
    \renewcommand{\appendixname}{Phụ lục}%
    \renewcommand{\indexname}{Chỉ mục}%
    \renewcommand{\figureautorefname}{Hình}%
    \renewcommand{\tableautorefname}{Bảng}%
    \renewcommand{\equationautorefname}{PT}%
    \renewcommand{\chapterautorefname}{Chương}%
    \renewcommand{\sectionautorefname}{mục}%
    \renewcommand{\subsectionautorefname}{mục}%
    \renewcommand{\appendixautorefname}{Phụ lục}%
}

% Khối chữ ký: "Hà Nội, ngày … tháng … năm …"
\newcommand{\khoichuky}[1]{%
    \par\vspace{12pt}
    \hfill
    \begin{minipage}{7cm}
        \centering\singlespacing
        \textit{Hà Nội, ngày …. tháng …. năm ….}\\
        \textbf{#1}\\
        \textit{Ký và ghi rõ họ tên}
    \end{minipage}\par
}

% Chèn trang scan (PDF) nếu đã có file, nếu chưa có thì in trang tiêu đề giữ chỗ
% #1: tên file PDF, #2: tiêu đề trang, #3: ghi chú giữ chỗ
\newcommand{\chentrangscan}[3]{%
    \IfFileExists{#1}{%
        \includepdf[pages=-]{#1}%
    }{%
        \chapter*{#2}%
        \textit{#3}%
        \clearpage
    }%
}

% Thông tin dùng chung
% Đánh dấu chỗ còn phải tự viết: in chữ đỏ trong ngoặc vuông.
% Trước khi nộp, tìm "\canviet" để chắc chắn không còn chỗ nào.
\newcommand{\canviet}[1]{\textcolor{red}{[\textit{#1}]}}

\newcommand{\tendetai}{Xây dựng hệ thống quản lý tín chỉ cacbon trên blockchain}
\newcommand{\thoigiannop}{1/2027}   % Tháng/năm nộp quyển, in ở cuối trang bìa

\usepackage[unicode,hidelinks,bookmarksnumbered]{hyperref}
% Cho phép ngắt dòng URL dài trong danh mục tài liệu tham khảo tại dấu gạch nối
\expandafter\def\expandafter\UrlBreaks\expandafter{\UrlBreaks\do\-}

\begin{document}

% ================= 1. BÌA TRƯỚC =================
\newgeometry{top=2cm,bottom=2cm,left=2.5cm,right=2cm}
\begin{titlepage}
    \setstretch{1}
    \setlength{\parskip}{0pt}
    \centering

    {\fontsize{14}{17}\selectfont\bfseries
        ĐẠI HỌC BÁCH KHOA HÀ NỘI\\
        KHOA TOÁN – TIN\par}

    \vspace{6.2cm}

    {\fontsize{18}{22}\selectfont\bfseries ĐỒ ÁN TỐT NGHIỆP\par}

    \vspace{22pt}

    % Tên đề tài phải khớp đúng với tên đã đăng ký trong phiếu giao nhiệm vụ
    {\fontsize{14}{17}\selectfont\bfseries \tendetai\par}

    \vspace{18pt}

    {\fontsize{14}{20}\selectfont
        \textbf{NGUYỄN CẢNH DUY}\\
        Email: Duy.NC227174@sis.hust.edu.vn\\
        Mã sinh viên: 20227174\\
        \textbf{Chuyên ngành Hệ thống thông tin quản lý}\par}

    \vspace{2.2cm}

    {\fontsize{13}{15.6}\selectfont
    \noindent\hspace*{0.6cm}%
    \makebox[5.3cm][l]{\textbf{Giảng viên hướng dẫn:}}%
    \makebox[5.8cm][l]{PGS.TS Lê Chí Ngọc}%
    \raisebox{-4pt}[0pt][0pt]{%
        \parbox[t]{3.5cm}{%
            \rule{3.5cm}{0.5pt}\\[-3pt]
            \hspace*{0.1cm}{\fontsize{10}{12}\selectfont Chữ ký của GVHD}}}%
    \par}

    \vfill

    {\fontsize{13}{15.6}\selectfont\bfseries HÀ NỘI, \thoigiannop\par}
\end{titlepage}
\restoregeometry

\pagenumbering{roman}

% ================= 2. NHẬN XÉT CỦA GIẢNG VIÊN HƯỚNG DẪN =================
% Bắt buộc theo Hướng dẫn thực hiện ĐA/TT của Khoa (02/04/2025): phải có chữ ký của GVHD.
% Khi có bản scan, đặt file NhanXetGVHD.pdf cùng thư mục, trang sẽ tự được chèn vào.
\chentrangscan{NhanXetGVHD.pdf}{Nhận xét của giảng viên hướng dẫn}
    {Chèn bản scan nhận xét có chữ ký của giảng viên hướng dẫn (file NhanXetGVHD.pdf).}

% ================= 3. PHIẾU BÁO CÁO TIẾN ĐỘ ĐỒ ÁN =================
% Phiếu in từ hệ thống số hóa của Khoa (fami.hust.edu.vn/sohoa), đóng vào quyển.
\chentrangscan{PhieuTienDo.pdf}{Phiếu báo cáo tiến độ đồ án}
    {Chèn phiếu báo cáo tiến độ in từ hệ thống số hóa (file PhieuTienDo.pdf).}

% ================= 4. LỜI CẢM ƠN + TÓM TẮT (CÙNG 1 TRANG, SINH VIÊN KÝ) =================
\clearpage
\begin{center}\vspace*{18pt}\textbf{Lời cảm ơn}\end{center}
% Tùy chọn, ngắn gọn, tránh từ sáo rỗng, khoảng 100–150 từ.
[Lời cảm ơn]

\begin{center}\vspace{12pt}\textbf{Tóm tắt nội dung đồ án}\end{center}
% Tối đa 300 chữ, cần nêu: vấn đề cần thực hiện; phương pháp thực hiện; công cụ sử dụng
% (phần mềm, phần cứng…); kết quả có phù hợp với vấn đề đặt ra không; tính thực tế;
% định hướng phát triển (nếu có); kiến thức và kỹ năng đã đạt được.
[Tóm tắt nội dung đồ án]

\khoichuky{Sinh viên thực hiện}

% ================= 5–7. MỤC LỤC & CÁC DANH MỤC =================
\tableofcontents

\clearpage
\phantomsection
\addcontentsline{toc}{chapter}{DANH MỤC HÌNH VẼ}
\listoffigures

\clearpage
\phantomsection
\addcontentsline{toc}{chapter}{DANH MỤC BẢNG BIỂU}
\listoftables

% Danh mục ký hiệu: bắt buộc theo hướng dẫn của Khoa. Liệt kê các chữ viết tắt theo thứ tự
% chữ cái, mỗi chữ viết tắt chỉ giải thích một lần, dùng thống nhất trong toàn bài.
\chapter*{Danh mục ký hiệu}
\phantomsection
\addcontentsline{toc}{chapter}{DANH MỤC KÝ HIỆU}
\begin{longtable}{@{}p{2.4cm}p{5.4cm}p{6.3cm}@{}}
\toprule
\textbf{Ký hiệu} & \textbf{Tiếng Anh} & \textbf{Ý nghĩa} \\
\midrule
\endhead
\bottomrule
\endfoot
AEF    & Agreed Electronic Format                  & Định dạng điện tử thỏa thuận \\
AFOLU  & Agriculture, Forestry and Other Land Use  & Nông nghiệp, lâm nghiệp và sử dụng đất khác \\
API    & Application Programming Interface         & Giao diện lập trình ứng dụng \\
CBAM   & Carbon Border Adjustment Mechanism        & Cơ chế điều chỉnh biên giới carbon \\
CDM    & Clean Development Mechanism               & Cơ chế phát triển sạch \\
CDOP   & Carbon Data Open Protocol                 & Giao thức mở dữ liệu carbon \\
CSDL   & Database                                  & Cơ sở dữ liệu \\
dApp   & Decentralized Application                 & Ứng dụng phi tập trung \\
DeFi   & Decentralized Finance                     & Tài chính phi tập trung \\
ERC    & Ethereum Request for Comments             & Chuẩn kỹ thuật của Ethereum \\
EVM    & Ethereum Virtual Machine                  & Máy ảo Ethereum \\
GHG    & Greenhouse Gas                            & Khí nhà kính \\
IoT    & Internet of Things                        & Internet vạn vật \\
IPFS   & InterPlanetary File System                & Hệ thống lưu trữ tệp phi tập trung \\
ITMO   & Internationally Transferred Mitigation Outcome & Kết quả giảm nhẹ được chuyển giao quốc tế \\
MRV    & Measurement, Reporting and Verification   & Đo đạc, báo cáo, thẩm định \\
NDC    & Nationally Determined Contribution        & Đóng góp do quốc gia tự quyết định \\
NFT    & Non-Fungible Token                        & Token không thể thay thế \\
OMGE   & Overall Mitigation in Global Emissions    & Giảm nhẹ tổng thể phát thải toàn cầu \\
PACM   & Paris Agreement Crediting Mechanism       & Cơ chế tín chỉ theo Điều 6.4 Thỏa thuận Paris \\
PDD    & Project Design Document                   & Tài liệu thiết kế dự án \\
PoS    & Proof of Stake                            & Bằng chứng cổ phần \\
PoW    & Proof of Work                             & Bằng chứng công việc \\
tCO\textsubscript{2}e & Tonne of CO\textsubscript{2} equivalent & Tấn CO\textsubscript{2} tương đương \\
VCS    & Verified Carbon Standard                  & Tiêu chuẩn carbon được xác minh của Verra \\
VCU    & Verified Carbon Unit                      & Đơn vị carbon được xác minh \\
VVB    & Validation and Verification Body          & Bên thẩm định độc lập \\
\end{longtable}

% ================= NỘI DUNG CHÍNH =================
% Theo hướng dẫn của Khoa: chương mở đầu → cơ sở lý thuyết → các chương kết quả,
% thực nghiệm → kết luận. Mỗi chương nên có đoạn mở đầu chương và đoạn tổng kết chương.
\clearpage
\pagenumbering{arabic}

% Chương 1 (Giai đoạn 1): khảo sát + hai mục Khoa yêu cầu là Nhiệm vụ đồ án và Cấu trúc đồ án.
% ============================================================
% CHƯƠNG 1 — GIỚI THIỆU VÀ KHẢO SÁT BÀI TOÁN  (Giai đoạn 1 của phiếu)
% Trích dẫn: \cite{...} từ References.bib
% Chỉ mục: \index{khoá sắp xếp@Tên hiển thị}
% Ghi chú: % [GÓP Ý] = gợi ý sửa, % [KIỂM TRA] = số liệu cần đối chiếu nguồn,
%          \canviet{...} = chỗ còn phải tự viết (in màu đỏ, xóa hết trước khi nộp).
% Các mục "Đổi mới số hóa MRV" và "Token hóa" đã chuyển sang Chương 2 (nghiên cứu liên quan),
% mục "Đề xuất kiến trúc tổng thể" đã chuyển sang Chương 3.
% ============================================================

\chapter{GIỚI THIỆU VÀ KHẢO SÁT BÀI TOÁN}
\label{chap:khaosat}
% ============================================================
% FIGURES_CHAP1.TEX — Hình tự vẽ (TikZ) cho Chương 1
% Mỗi hình là một lệnh, gọi đúng chỗ trong Chap1.tex:
%   \HinhDieuSau  \HinhMRV  \HinhVaccargiu  \HinhToucan  \HinhKienTruc
% Các hình được vẽ vừa bề rộng vùng chữ (15 cm), không cần thu nhỏ.
% Preamble main.tex cần:
%   \usepackage{tikz}
%   \usetikzlibrary{shapes.geometric,arrows.meta,positioning,calc}
% ============================================================

\tikzset{
    hop/.style      = {draw, rounded corners=3pt, align=center,
                       inner sep=3pt, font=\small},
    muiten/.style   = {-{Stealth[length=6pt]}, thick},
    ghichu/.style   = {font=\footnotesize\itshape, text=black!65, align=center},
}

% ------------------------------------------------------------
% HÌNH: Điều 6.2 và Điều 6.4 — đặt sau Bảng so sánh Điều 6
% ------------------------------------------------------------
\newcommand{\HinhDieuSau}{%
\begin{figure}[htbp]
\centering
\begin{tikzpicture}
    % Tiêu đề hai cột và đường phân cách
    \node[font=\bfseries] at (3.4, 0) {Điều 6.2 (song phương)};
    \node[font=\bfseries] at (11.0, 0) {Điều 6.4 (tập trung – PACM)};
    \draw[thick, black!30] (7.3, 0.3) -- (7.3, -7.0);

    % ---- Điều 6.2 ----
    \node[hop, fill=blue!10, minimum width=3.0cm, minimum height=1.1cm]
        (A) at (1.7, -1.3) {\textbf{Quốc gia A}\\(bán ITMOs)};
    \node[hop, fill=blue!10, minimum width=3.0cm, minimum height=1.1cm]
        (B) at (1.7, -4.1) {\textbf{Quốc gia B}\\(mua ITMOs)};
    \node[draw, diamond, aspect=1.8, fill=yellow!25, font=\small,
          inner sep=1pt] (I) at (5.3, -2.8) {ITMOs};
    \draw[muiten] (A.east) -| node[ghichu, pos=0.25, above] {bán} (I.north);
    \draw[muiten] (I.south) |- node[ghichu, pos=0.75, above] {chuyển giao} (B.east);
    \node[hop, fill=red!8, text width=6.2cm, font=\footnotesize]
        (CA) at (3.4, -6.4)
        {\textbf{Điều chỉnh tương ứng:}\\
         Quốc gia A cộng thêm phát thải vào sổ quốc gia;\\
         Quốc gia B được trừ lượng phát thải tương ứng};
    \draw[muiten, dashed, red!60] (B.south) -- (B.south |- CA.north);

    % ---- Điều 6.4 ----
    \node[hop, fill=green!12, minimum width=3.2cm, minimum height=0.9cm]
        (P) at (10.2, -1.2) {Nhà phát triển dự án};
    \node[hop, fill=orange!15, minimum width=3.2cm, minimum height=0.9cm]
        (S) at (10.2, -2.7) {Cơ quan Giám sát};
    \node[draw, diamond, aspect=1.8, fill=yellow!25, font=\small,
          inner sep=1pt] (E) at (10.2, -4.3) {A6.4ERs};
    \node[hop, fill=green!12, minimum width=3.2cm, minimum height=0.9cm]
        (U) at (10.2, -5.9) {Người mua / bù trừ};
    \draw[muiten] (P) -- node[ghichu, right] {đăng ký} (S);
    \draw[muiten] (S) -- node[ghichu, left] {cấp} (E);
    \draw[muiten] (E) -- node[ghichu, right] {giao dịch} (U);
    % Khoản trích 5%/2% áp dụng tại thời điểm cấp (quyết định 3/CMA.3), không phải mỗi giao dịch
    \node[hop, fill=red!8, text width=2.3cm, font=\footnotesize]
        (F) at (13.7, -3.5) {Trích khi cấp:\\5\% Quỹ Thích ứng\\$\geq$2\% hủy (OMGE)};
    \coordinate (SE) at ($(S.south)!0.5!(E.north)$);
    \draw[muiten, red!60] (SE) -- (SE -| F.west);
\end{tikzpicture}
\caption[Cơ chế hoạt động của Điều 6.2 và Điều 6.4 Thỏa thuận Paris]%
{Cơ chế hoạt động của Điều 6.2 và Điều 6.4 Thỏa thuận Paris~\cite{columbia2025article6} \cite{unfccc2022cma3}.}
\label{fig:dieu6}
\end{figure}}

% ------------------------------------------------------------
% HÌNH: Vòng đời tín chỉ và ISO 14064 — đặt sau Bảng ISO 14064
% ------------------------------------------------------------
\newcommand{\HinhMRV}{%
\begin{figure}[htbp]
\centering
\begin{tikzpicture}[
    buoc/.style = {draw, rounded corners=4pt, align=center,
                   text width=3.4cm, minimum height=1.5cm, font=\small},
]
    % Hàng 1: trái -> phải
    \node[buoc, fill=blue!12]   (s1) at (2.0, 0)
        {\textbf{1. Thiết kế dự án}\\{\footnotesize\itshape PDD – ISO 14064-2}};
    \node[buoc, fill=teal!15]   (s2) at (7.4, 0)
        {\textbf{2. Thẩm định}\\{\footnotesize\itshape VVB – ISO 14064-3}};
    \node[buoc, fill=green!15]  (s3) at (12.8, 0)
        {\textbf{3. Giám sát}\\{\footnotesize\itshape dữ liệu hoạt động}};
    % Hàng 2: phải -> trái
    \node[buoc, fill=yellow!22] (s4) at (12.8, -2.6)
        {\textbf{4. Xác minh}\\{\footnotesize\itshape VVB – ISO 14064-3}};
    \node[buoc, fill=orange!20] (s5) at (7.4, -2.6)
        {\textbf{5. Phát hành}\\{\footnotesize\itshape trừ tài khoản đệm}};
    \node[buoc, fill=red!12]    (s6) at (2.0, -2.6)
        {\textbf{6. Giao dịch \& hủy}\\{\footnotesize\itshape hủy vĩnh viễn (retire)}};

    \draw[muiten] (s1) -- (s2);
    \draw[muiten] (s2) -- (s3);
    \draw[muiten] (s3) -- (s4);
    \draw[muiten] (s4) -- (s5);
    \draw[muiten] (s5) -- (s6);
\end{tikzpicture}
\caption[Vòng đời tín chỉ carbon và các phần tương ứng của tiêu chuẩn ISO 14064]%
{Vòng đời tín chỉ carbon và các phần tương ứng của tiêu chuẩn ISO 14064~\cite{iso14064-2} \cite{iso14064-3}.}
\label{fig:mrv}
\end{figure}}

% ------------------------------------------------------------
% HÌNH: Kiến trúc 4 lớp của Vaccargiu — đặt sau danh sách 4 lớp (Chương 2, mục \ref{sec:dmrv})
% ------------------------------------------------------------
\newcommand{\HinhVaccargiu}{%
\begin{figure}[htbp]
\centering
\begin{tikzpicture}[
    lop/.style = {draw, rounded corners=5pt, align=center,
                  text width=9.0cm, minimum height=1.3cm, font=\small},
]
    \node[lop, fill=orange!15] (l4) at (4.8, 5.4)
        {\textbf{Lớp 4: Chuỗi khối và chứng nhận}\\
         {\footnotesize Hyperledger Fabric, hợp đồng thông minh tạo và đăng ký tín chỉ}};
    \node[lop, fill=yellow!18] (l3) at (4.8, 3.6)
        {\textbf{Lớp 3: Tổng hợp dữ liệu tại biên}\\
         {\footnotesize tổng hợp theo phút, gom thành lô 5 phút để ghi lên chuỗi}};
    \node[lop, fill=blue!10]   (l2) at (4.8, 1.8)
        {\textbf{Lớp 2: Thu thập dữ liệu}\\
         {\footnotesize cổng nối nhận luồng dữ liệu thời gian thực}};
    \node[lop, fill=green!13]  (l1) at (4.8, 0)
        {\textbf{Lớp 1: Cảm biến vật lý}\\
         {\footnotesize 8 đồng hồ đo thông minh 3 pha (24 pha điện)}};

    \draw[muiten] (l1) -- (l2);
    \draw[muiten] (l2) -- (l3);
    \draw[muiten] (l3) -- (l4);

    % Chú thích bên phải: tần suất / dạng dữ liệu
    \node[ghichu, anchor=west, text width=4.3cm] at (10.2, 5.4) {288 lần ghi / ngày};
    \node[ghichu, anchor=west, text width=4.3cm] at (10.2, 3.6) {1.440 bản ghi / ngày};
    \node[ghichu, anchor=west, text width=4.3cm] at (10.2, 1.8) {luồng dữ liệu liên tục};
    \node[ghichu, anchor=west, text width=4.3cm] at (10.2, 0)   {1–2 giây / lần đo};
\end{tikzpicture}
\caption[Kiến trúc 4 lớp của hệ thống chứng nhận tín chỉ carbon dựa trên IoT và blockchain]%
{Kiến trúc 4 lớp của hệ thống chứng nhận tín chỉ carbon dựa trên IoT và blockchain~\cite{vaccargiu2026blockchain}.}
\label{fig:vaccargiu}
\end{figure}}

% ------------------------------------------------------------
% HÌNH: Luồng token Toucan — đặt sau danh sách các bước Toucan (Chương 2, mục \ref{sec:token})
% Theo tài liệu Toucan: Batch NFT là ERC-721, TCO2 là ERC-20 (mỗi dự án–vintage một hợp đồng),
% retire() đốt TCO2; cầu nối Verra đã ngừng hoạt động.
% ------------------------------------------------------------
\newcommand{\HinhToucan}{%
\begin{figure}[htbp]
\centering
\begin{tikzpicture}[
    tok/.style = {draw, rounded corners=4pt, align=center,
                  text width=2.2cm, minimum height=1.3cm, font=\small},
]
    \node[tok, fill=blue!12]   (v) at (1.2, 0)    {\textbf{Verra}\\VCUs};
    \node[tok, fill=teal!15]   (b) at (4.9, 0)    {\textbf{Batch NFT}\\ERC-721};
    \node[tok, fill=teal!25]   (t) at (8.6, 0)    {\textbf{TCO2}\\ERC-20};
    \node[tok, fill=green!15]  (p) at (12.6, 0)   {\textbf{BCT / NCT}\\ERC-20};
    \node[tok, fill=red!12]    (r) at (8.6, -2.8) {\textbf{Hủy}\\burn TCO2};
    \node[tok, fill=yellow!20, text width=2.6cm] (c) at (12.6, -2.8) {\textbf{Chứng nhận}\\hủy (NFT)};

    \draw[muiten] (v) -- node[above, font=\scriptsize] {cầu nối}  (b);
    \draw[muiten] (b) -- node[above, font=\scriptsize] {chia nhỏ} (t);
    \draw[muiten] ([yshift=3mm]t.east) -- node[above, font=\scriptsize\ttfamily] {deposit} ([yshift=3mm]p.west);
    \draw[muiten] ([yshift=-3mm]p.west) -- node[below, font=\scriptsize\ttfamily] {redeem} ([yshift=-3mm]t.east);
    \draw[muiten] (t) -- node[left, font=\scriptsize\ttfamily] {retire()} (r);
    \draw[muiten] (r) -- (c);

    \node[ghichu, text width=2.8cm] at (1.2, -1.5) {VCU được hủy trên\\Verra, ghi mã lô};
    \node[ghichu, text width=2.8cm] at (4.9, -1.5) {Toucan duyệt\\số sê-ri};
    \node[ghichu, text width=3.0cm] at (12.6, -1.5) {1 token = 1 tCO$_2$e};
\end{tikzpicture}
\caption[Luồng token hóa và hủy tín chỉ carbon qua Toucan Protocol]%
{Luồng token hóa và hủy tín chỉ carbon qua Toucan Protocol (cầu nối Verra hiện đã ngừng hoạt động)~\cite{toucan2025docs}.}
\label{fig:toucan}
\end{figure}}

% ------------------------------------------------------------
% HÌNH: Kiến trúc tổng thể đề xuất — đặt cuối mục \ref{sec:kientruc} (Chương 3)
% Đây là thiết kế của đồ án nên không trích nguồn. Ô nét đứt = ngoài phạm vi cài đặt.
% ------------------------------------------------------------
\newcommand{\HinhKienTruc}{%
\begin{figure}[htbp]
\centering
\begin{tikzpicture}[
    tp/.style = {draw, rounded corners=3pt, align=center, fill=white,
                 text width=3.0cm, minimum height=1.0cm, font=\footnotesize},
    lp/.style = {draw, very thick, rounded corners=6pt,
                 minimum width=14.6cm, minimum height=2.1cm},
    tl/.style = {font=\small\bfseries, anchor=north west, xshift=6pt, yshift=-3pt},
]
    % Lớp 4
    \node[lp, fill=purple!10] (L4) at (7.3, 0) {};
    \node[tl] at (L4.north west) {Lớp 4 — Giao diện ứng dụng \& API};
    \node[tp] at (1.9, -0.3) {Next.js, ethers.js\\giao diện web};
    \node[tp] at (5.5, -0.3) {MetaMask\\ký giao dịch};
    \node[tp] at (9.1, -0.3) {API REST\\(CAD Trust V2.0)};
    \node[tp] at (12.7,-0.3) {Dashboard\\tra cứu, thống kê};

    % Lớp 3
    \node[lp, fill=orange!12] (L3) at (7.3, -2.7) {};
    \node[tl] at (L3.north west) {Lớp 3 — Hợp đồng thông minh (on-chain, Ethereum Sepolia)};
    \node[tp] at (1.9, -3.0) {AccessControl\\VERIFIER\_ROLE};
    \node[tp] at (5.5, -3.0) {Minting\\NFT ERC-721};
    \node[tp] at (9.1, -3.0) {Pool\\token ERC-20};
    \node[tp] at (12.7,-3.0) {Retirement\\burn + sự kiện};

    % Lớp 2
    \node[lp, fill=blue!10] (L2) at (7.3, -5.4) {};
    \node[tl] at (L2.north west) {Lớp 2 — Dịch vụ \& lưu trữ off-chain};
    \node[tp] at (1.9, -5.7) {Backend\\Node.js, Express};
    \node[tp] at (5.5, -5.7) {PostgreSQL\\(CAD Trust V2.0)};
    \node[tp] at (9.1, -5.7) {IPFS / S3\\PDD, metadata};
    \node[tp] at (12.7,-5.7) {Event listener\\đồng bộ sự kiện};

    % Lớp 1
    \node[lp, fill=green!12] (L1) at (7.3, -8.1) {};
    \node[tl] at (L1.north west) {Lớp 1 — Nguồn dữ liệu \& MRV};
    \node[tp] at (1.9, -8.4) {Chủ dự án\\hồ sơ, PDD};
    \node[tp] at (5.5, -8.4) {VVB\\báo cáo thẩm định};
    \node[tp] at (9.1, -8.4) {Dữ liệu giám sát\\(mô phỏng)};
    \node[tp, dashed] at (12.7,-8.4) {IoT, dMRV\\(hướng phát triển)};

    % Liên kết giữa các lớp
    \foreach \tren/\duoi in {L4/L3, L3/L2, L2/L1}
        \draw[{Stealth[length=6pt]}-{Stealth[length=6pt]}, very thick, black!45]
            (\duoi.north) -- (\tren.south);
\end{tikzpicture}
\caption[Kiến trúc tổng thể 4 lớp của hệ thống quản lý tín chỉ carbon đề xuất]%
{Kiến trúc tổng thể 4 lớp của hệ thống quản lý tín chỉ carbon đề xuất (mô hình hỗn hợp on-chain/off-chain).}
\label{fig:kientruc}
\end{figure}}
  % Tất cả hình TikZ – xem file figures_chap1.tex

% Đoạn mở đầu chương (3–5 câu): chương này khảo sát những gì (thị trường, tiêu chuẩn MRV,
% hệ thống đăng ký, pháp lý Việt Nam) và dẫn tới nhiệm vụ của đồ án ở cuối chương.
\canviet{Đoạn mở đầu chương}

% ------------------------------------------------------------
\section{Động lực thị trường và khung pháp lý khí hậu toàn cầu}
\label{sec:dongluc}

Biến đổi khí hậu là một trong những thách thức lớn hiện nay; vì vậy, việc chuyển sang các nguồn năng lượng tái tạo và giảm phát thải khí nhà kính (GHG)\index{GHG|see{Khí nhà kính}}\index{Khi nha kinh@Khí nhà kính} ngày càng trở thành một yêu cầu pháp lý và kinh tế, không chỉ là một mối quan tâm về môi trường~\cite{liu2024blockchain} \cite{columbia2025article6}. Thị trường carbon\index{Thi truong carbon@Thị trường carbon} toàn cầu, gồm thị trường tuân thủ\index{Thi truong carbon@Thị trường carbon!tuan thu@tuân thủ} và thị trường tự nguyện\index{Thi truong carbon@Thị trường carbon!tu nguyen@tự nguyện}, là cơ chế tài chính dùng để định giá lượng phát thải carbon. Hai thị trường này có quy mô rất khác nhau. Năm 2023, các thị trường hạn ngạch phát thải (thị trường tuân thủ) đạt giá trị kỷ lục khoảng 881 tỷ EUR (khoảng 949 tỷ USD) với khoảng 12,5 tỷ tấn hạn ngạch được giao dịch, trong đó hệ thống EU ETS chiếm khoảng 87\% giá trị~\cite{lseg2024yir}. Trong cùng năm, thị trường tín chỉ tự nguyện chỉ đạt khoảng 723 triệu USD với khoảng 111 triệu tấn CO$_2$e được giao dịch~\cite{foresttrends2024sovcm}. Một tín chỉ carbon\index{Tin chi carbon@Tín chỉ carbon} được tiêu chuẩn hóa đại diện cho một tấn carbon dioxide tương đương (tCO$_2$e)\index{tCO2e@tCO$_2$e} đã được tránh, giảm hoặc loại bỏ khỏi khí quyển thông qua các hoạt động được thẩm định độc lập~\cite{verra2024vcs} \cite{columbia2025article6}.

Dù đã đạt quy mô lớn, thị trường tín chỉ carbon vẫn còn các hạn chế về tính minh bạch, khả năng tiếp cận và đặc biệt là rủi ro tính toán kép (double counting)\index{Tinh toan kep@Tính toán kép (double counting)}~\cite{schneider2019double}. Theo Schneider và cộng sự, tính toán kép có ba dạng chính~\cite{schneider2019double}: phát hành kép (double issuance -- nhiều tín chỉ được cấp cho cùng một lượng giảm phát thải, ví dụ một dự án đăng ký ở hai chương trình khác nhau)\index{Tinh toan kep@Tính toán kép (double counting)!phat hanh kep@phát hành kép}, sử dụng kép (double use -- cùng một tín chỉ được bán hoặc sử dụng hai lần)\index{Tinh toan kep@Tính toán kép (double counting)!su dung kep@sử dụng kép} và tuyên bố kép (double claiming -- cả quốc gia sở tại và bên mua đều tính cùng một lượng giảm phát thải vào mục tiêu của mình)\index{Tinh toan kep@Tính toán kép (double counting)!tuyen bo kep@tuyên bố kép}. Để xử lý các vấn đề này ở cấp quốc gia, Thỏa thuận Paris (Paris Agreement)\index{Thoa thuan Paris@Thỏa thuận Paris} đưa ra Điều~6\index{Thoa thuan Paris@Thỏa thuận Paris!Dieu 6@Điều 6}, quy định khung hạch toán và các cơ chế hợp tác về thị trường carbon giữa các quốc gia. Do đó, cần nắm được cấu trúc của Điều~6 trước khi thiết kế một hệ thống phần mềm quản lý tín chỉ carbon.

Điều~6 gồm hai cơ chế chính với cách vận hành và yêu cầu dữ liệu khác nhau. Khía cạnh đầu tiên, Điều~6.2\index{Thoa thuan Paris@Thỏa thuận Paris!Dieu 6.2@Điều 6.2}, thiết lập một khuôn khổ phi tập trung cho các thỏa thuận hợp tác song phương hoặc đa phương trực tiếp giữa các quốc gia~\cite{columbia2025article6}. Thông qua Điều~6.2, các quốc gia có thể giao dịch các Kết quả giảm nhẹ được chuyển giao quốc tế (Internationally Transferred Mitigation Outcomes -- ITMOs)\index{ITMOs} để đạt được Đóng góp do quốc gia tự quyết định (NDCs)\index{NDCs}~\cite{columbia2025article6}. Đặc điểm kỹ thuật quan trọng của Điều~6.2 là yêu cầu thực hiện ``Điều chỉnh tương ứng'' (Corresponding Adjustments)\index{Dieu chinh tuong ung@Điều chỉnh tương ứng (Corresponding Adjustments)}. Về mặt hạch toán, khi một quốc gia chuyển giao ITMOs, quốc gia đó phải cộng lượng phát thải tương đương vào sổ hạch toán khí nhà kính quốc gia của mình, còn quốc gia nhận được trừ đi lượng tương ứng; nhờ đó cùng một kết quả giảm phát thải không được cả hai bên tính vào NDC (tránh tuyên bố kép)~\cite{columbia2025article6} \cite{schneider2019double}. Giao dịch chuyển giao ITMO đầu tiên theo Điều~6.2 được hoàn tất giữa Thái Lan và Thụy Sĩ: 1.916 ITMO từ Chương trình Xe buýt điện Bangkok được ghi có vào hệ thống đăng ký của Thụy Sĩ ngày 15/12/2023 và được công bố vào tháng 1 năm 2024, bên mua là Quỹ KliK của Thụy Sĩ~\cite{spglobal2024itmo} \cite{columbia2025article6}. Các ITMO có điều chỉnh tương ứng thường có chi phí cao hơn tín chỉ tự nguyện. Quỹ KliK dự kiến chi phí trung bình khoảng 35~CHF/tấn CO$_2$ trong giai đoạn 2022--2024 và 27~CHF/tấn trong giai đoạn 2025--2030~\cite{argus2024itmo}, trong khi giá trung bình của tín chỉ trên thị trường tự nguyện năm 2023 là khoảng 6,53~USD/tấn~\cite{foresttrends2024sovcm}. Một phần chênh lệch này thường được giải thích là ``phần bù tính toàn vẹn'' (integrity premium)\index{Phan bu tinh toan ven@Phần bù tính toàn vẹn (integrity premium)}, tuy nhiên hai mức giá không so sánh trực tiếp được vì khác nhau về loại dự án, điều kiện hợp đồng và chi phí của bên mua.

Cơ chế thứ hai, Điều~6.4\index{Thoa thuan Paris@Thỏa thuận Paris!Dieu 6.4@Điều 6.4}, là Cơ chế cấp tín chỉ của Thỏa thuận Paris (PACM)\index{PACM}, kế thừa Cơ chế Phát triển Sạch (CDM)\index{CDM} của Nghị định thư Kyoto\index{Nghi dinh thu Kyoto@Nghị định thư Kyoto}~\cite{columbia2025article6} \cite{unfccc2022cma3}. Khác với cách tiếp cận phi tập trung của Điều~6.2, Điều~6.4 là một cơ chế tập trung do Liên Hợp Quốc quản lý thông qua một Cơ quan Giám sát (Supervisory Body) gồm 12 thành viên~\cite{columbia2025article6}. Ngày 26/02/2026, Cơ quan Giám sát phê duyệt đợt cấp tín chỉ đầu tiên của PACM cho một chương trình bếp đun cải tiến tại Myanmar, vốn được chuyển tiếp từ CDM và phối hợp với các bên tham gia được Hàn Quốc ủy quyền~\cite{unfccc2026pacm}. Hai yêu cầu cấp được duyệt có tổng cộng 648.783 tín chỉ cho giai đoạn giám sát từ tháng 1/2021 đến tháng 5/2022. Lượng này thấp hơn khoảng 40\% so với ước tính theo báo cáo giám sát của CDM, chủ yếu do tỷ lệ sinh khối không tái tạo (fNRB) được điều chỉnh từ 0,615 xuống giá trị mặc định 0,36 của Myanmar~\cite{unfccc2026pacm} \cite{fastmarkets2026pacm}. Ngoài ra, hệ thống dữ liệu cho Điều~6.4 phải xử lý các khoản trích bắt buộc tại thời điểm cấp tín chỉ (không phải ở mỗi giao dịch): 5\% lượng A6.4ER được chuyển vào Quỹ Thích ứng (Adaptation Fund)\index{Quy Thich ung@Quỹ Thích ứng (Adaptation Fund)} để hỗ trợ các quốc gia đang phát triển, và tối thiểu 2\% bị hủy để góp phần giảm nhẹ tổng thể lượng phát thải toàn cầu (Overall Mitigation in Global Emissions -- OMGE)\index{OMGE}~\cite{unfccc2022cma3}. Bảng~\ref{tab:dieu6} tổng hợp sự khác biệt giữa hai cơ chế.

\begin{table}[htbp]
    \centering
    \caption{So sánh yêu cầu hệ thống giữa Điều 6.2 và Điều 6.4 của Thỏa thuận Paris}
    \label{tab:dieu6}
    \small
    \renewcommand{\arraystretch}{1.25}
    \begin{tabularx}{\textwidth}{|p{3.2cm}|X|X|}
        \hline
        \textbf{Yêu cầu / Tiêu chí} & \textbf{Điều 6.2 (Thỏa thuận song phương)} & \textbf{Điều 6.4 (Cơ chế tập trung của Liên Hợp Quốc)} \\
        \hline
        Bản chất thị trường & Phi tập trung, đàm phán trực tiếp giữa chính phủ các quốc gia~\cite{columbia2025article6}. & Tập trung, được giám sát bởi Cơ quan Giám sát của Liên Hợp Quốc~\cite{columbia2025article6}. \\
        \hline
        Đơn vị giao dịch & Kết quả giảm nhẹ phát thải quốc tế (ITMOs)~\cite{columbia2025article6}. & Đơn vị Giảm phát thải Điều 6.4 (A6.4ERs)\index{A6.4ERs}~\cite{columbia2025article6}. \\
        \hline
        Điều chỉnh tương ứng & Bắt buộc đối với ITMO được ủy quyền và chuyển giao quốc tế~\cite{columbia2025article6} \cite{unfccc2022cma3}. & Chỉ áp dụng cho A6.4ER được quốc gia sở tại ủy quyền dùng cho NDC hoặc mục đích quốc tế khác~\cite{unfccc2022cma3}. \\
        \hline
        Khoản trích bắt buộc & Khuyến khích đóng góp (không bắt buộc)~\cite{unfccc2022cma3}. & Bắt buộc tại thời điểm cấp: 5\% cho Quỹ Thích ứng, tối thiểu 2\% hủy cho OMGE~\cite{unfccc2022cma3}. \\
        \hline
        Đối tượng tham gia & Chủ yếu là cấp Chính phủ quốc gia~\cite{columbia2025article6}. & Các nhà phát triển dự án toàn cầu, chính phủ và khu vực tư nhân~\cite{columbia2025article6}. \\
        \hline
    \end{tabularx}
\end{table}

\HinhDieuSau

\autoref{fig:dieu6} minh họa luồng hoạt động của hai cơ chế trên.

Sự khác biệt giữa hai cơ chế này đặt ra các yêu cầu khác nhau đối với hệ thống đăng ký và phần mềm quản lý tín chỉ. Hệ thống cần xử lý được các kịch bản hạch toán khác nhau, quản lý trạng thái của ``Thư ủy quyền'' (Letter of Authorization)\index{Thu uy quyen@Thư ủy quyền (Letter of Authorization)} từ quốc gia sở tại và ghi nhận đúng các khoản trích theo tỷ lệ phần trăm trong vòng đời của một tín chỉ.

% ------------------------------------------------------------
\section{Tiêu chuẩn hạch toán và quy trình đo đạc, báo cáo, thẩm định (MRV)}
\label{sec:mrv}

Để tín chỉ carbon được công nhận và giao dịch trên toàn cầu, toàn bộ vòng đời của chúng từ giai đoạn khởi tạo dự án, giám sát dữ liệu đến cấp phát phải tuân thủ các khung tiêu chuẩn độc lập. Để thiết kế một phần mềm MRV (đo đạc, báo cáo, thẩm định)\index{MRV}, cần nắm được cấu trúc dữ liệu được quy định bởi các tổ chức thiết lập tiêu chuẩn (Standard Setters) như Tổ chức Tiêu chuẩn hóa Quốc tế (ISO) và Verra\index{Verra}.

Bộ tiêu chuẩn ISO~14064\index{ISO 14064} được sử dụng rộng rãi để định lượng, báo cáo và thẩm định phát thải khí nhà kính~\cite{iso14064-1} \cite{iso14064-2} \cite{iso14064-3}. Bộ tiêu chuẩn gồm ba phần (Bảng~\ref{tab:iso14064}), tương ứng với các nhóm chức năng mà một phần mềm quản lý phát thải cần hỗ trợ.

\begin{table}[htbp]
    \centering
    \caption{Cấu trúc ba phần của tiêu chuẩn ISO 14064}
    \label{tab:iso14064}
    \small
    \renewcommand{\arraystretch}{1.25}
    \begin{tabularx}{\textwidth}{|p{2.3cm}|X|X|}
        \hline
        \textbf{Phần tiêu chuẩn} & \textbf{Phạm vi và chức năng cốt lõi} & \textbf{Phương pháp tiếp cận đo lường} \\
        \hline
        ISO 14064-1\index{ISO 14064!Phan 1@Phần 1 (cấp tổ chức)} & Áp dụng ở cấp độ tổ chức. Quản lý hệ thống kiểm kê GHG tổng thể~\cite{iso14064-1}. & Tính toán và phân loại phát thải theo phạm vi (Scope 1, Scope 2, Scope 3)\index{Scope 1, 2, 3} bằng cách nhân dữ liệu hoạt động với hệ số phát thải (Emission Factors)\index{He so phat thai@Hệ số phát thải (Emission Factor)}~\cite{iso14064-1}. \\
        \hline
        ISO 14064-2\index{ISO 14064!Phan 2@Phần 2 (cấp dự án)} & Áp dụng ở cấp độ dự án (ví dụ: dự án năng lượng tái tạo, trồng rừng)~\cite{iso14064-2}. & Lượng hóa lượng GHG được giảm hoặc loại bỏ bằng cách so sánh lượng phát thải thực tế với đường phát thải cơ sở (baseline)~\cite{iso14064-2}. \\
        \hline
        ISO 14064-3\index{ISO 14064!Phan 3@Phần 3 (validation, verification)} & Khung validation và verification~\cite{iso14064-3}. & Thiết lập tiêu chí và quy trình cho các bên thẩm định độc lập (VVB)\index{VVB} nhằm đảm bảo dữ liệu minh bạch, toàn vẹn và không có sai sót vật chất~\cite{iso14064-3}. \\
        \hline
    \end{tabularx}
\end{table}

\HinhMRV

Vòng đời tín chỉ carbon và sự tương ứng của từng phần tiêu chuẩn được thể hiện trong \autoref{fig:mrv}.

Trong đó, ISO~14064-2 định hình trực tiếp logic thuật toán của quá trình sinh tín chỉ. Nó yêu cầu hệ thống phần mềm phải lưu trữ được các ``đường phát thải cơ sở''\index{Duong phat thai co so@Đường phát thải cơ sở (baseline)} (đại diện cho mức phát thải nếu dự án không diễn ra) và chứng minh được ``Tính bổ sung'' (Additionality)\index{Tinh bo sung@Tính bổ sung (Additionality)} -- tức là lượng giảm phát thải sẽ không xảy ra nếu không có dự án và nguồn thu từ tín chỉ carbon, vượt ra ngoài kịch bản hoạt động thông thường (Business-as-Usual)~\cite{columbia2025article6} \cite{iso14064-2}.

Bên cạnh các tiêu chuẩn định lượng chung, Chương trình Tiêu chuẩn Carbon được Xác minh (Verified Carbon Standard -- VCS)\index{VCS (Verified Carbon Standard)} do Verra quản lý là chương trình cấp tín chỉ GHG được sử dụng rộng rãi nhất theo công bố của Verra. Tại thời điểm truy cập, chương trình có hơn 2.579 dự án tại hơn 132 quốc gia và đã cấp tín chỉ cho hơn 1,3 tỷ tấn GHG được giảm hoặc loại bỏ~\cite{verra2024vcs}. Dữ liệu của một dự án trên hệ thống Verra gồm nhiều thành phần và phải tuân theo các yêu cầu phương pháp luận (Methodology Requirements) riêng cho từng lĩnh vực, chẳng hạn Nông nghiệp, Lâm nghiệp và Sử dụng Đất (AFOLU)\index{AFOLU}, các chất làm suy giảm tầng ô-dôn (ODS) hoặc Lưu trữ Carbon Địa chất (GCS)~\cite{verra2024vcs}. Ngoài yêu cầu của chương trình, chất lượng tín chỉ còn được các tổ chức xếp hạng độc lập đánh giá theo các tiêu chí như tính bổ sung, tính lâu dài (permanence), độ tin cậy của định lượng, rò rỉ phát thải (leakage)\index{Ro ri phat thai@Rò rỉ phát thải (leakage)} và mức đáp ứng Nguyên tắc Carbon Cốt lõi (Core Carbon Principles) của ICVCM\index{ICVCM}. Vì vậy, hệ thống quản lý nên có trường dữ liệu lưu các kết quả xếp hạng này; mô hình dữ liệu CAD Trust V2.0 cũng có bảng Rating cho mục đích tương tự~\cite{cadtrust2025datamodel}. Đối với rủi ro đảo ngược (lượng carbon đã lưu giữ bị phát thải trở lại, ví dụ do cháy rừng), Verra duy trì tài khoản đệm gộp (pooled buffer account)\index{Tai khoan dem gop@Tài khoản đệm gộp (pooled buffer account)}: một tỷ lệ tín chỉ của dự án AFOLU, xác định bằng công cụ đánh giá rủi ro không vĩnh viễn, được giữ lại như một hình thức bảo hiểm chung cho toàn chương trình~\cite{verra2024vcs} \cite{verra2024projecthub}. Một phần mềm tích hợp với Verra cần tính đến khoản giữ lại này trong mô-đun phát hành (Issuance module) trước khi ghi Đơn vị carbon được xác minh (VCU)\index{VCU} vào tài khoản của nhà phát triển dự án.

% ------------------------------------------------------------
\section{Kiến trúc hệ thống đăng ký tập trung và nền tảng liên thông dữ liệu}
\label{sec:sodangky}

% [GÓP Ý] Đã bỏ câu ``Khảo sát hệ thống (system survey) là một yêu cầu bắt buộc để xây dựng báo cáo
% đồ án phần mềm'': câu này nói về yêu cầu viết báo cáo, không phải nội dung khảo sát.
Phần lớn hệ thống đăng ký\index{He thong dang ky@Hệ thống đăng ký (registry)} hiện nay hoạt động biệt lập trên cơ sở dữ liệu tập trung. Một số sáng kiến gần đây hướng tới liên thông dữ liệu giữa các hệ thống đăng ký, trong đó có sáng kiến dùng chuỗi khối, nhằm giảm tình trạng phân mảnh thông tin, giảm chi phí quản trị và hỗ trợ phát hiện tính toán kép giữa các hệ thống.

\subsection*{Hệ thống đăng ký tập trung: Verra Registry và UNDP National Carbon Registry}
\label{subsec:verra-undp}

Verra Registry\index{He thong dang ky@Hệ thống đăng ký (registry)!Verra Registry} đóng vai trò là điểm neo tập trung cho các dự án và tín chỉ VCUs. Nền tảng này theo dõi toàn bộ vòng đời của tín chỉ thông qua các quy trình được kiểm soát chặt chẽ: mở tài khoản, nộp hồ sơ yêu cầu phát hành, tạo số sê-ri duy nhất, và cơ chế hủy tín chỉ (retirement/cancellation)\index{Huy tin chi@Hủy tín chỉ (retirement)}~\cite{verra2024vcs}. Gần đây, Verra đưa vào sử dụng Verra Project Hub\index{Verra Project Hub}, nền tảng trực tuyến để nộp và xét duyệt hồ sơ dự án: công cụ theo dõi dự án (Project Tracker) ra mắt tháng 8/2023 và công cụ nộp hồ sơ số (Digital Project Submission Tool) ra mắt tháng 8/2024~\cite{verra2024projecthub}. Project Hub có các phương pháp luận được số hóa (ví dụ VM0038, VM0042) để tính lượng giảm và loại bỏ phát thải (ERR) từ dữ liệu nhập vào biểu mẫu, cùng các công cụ đánh giá rủi ro không vĩnh viễn (Non-Permanence Risk Tool) dạng số cho dự án AFOLU và GCS~\cite{verra2024projecthub}. Tuy vậy, tại thời điểm khảo sát, Verra chưa công bố tài liệu API công khai cho Verra Registry; dữ liệu công khai chủ yếu được tra cứu qua giao diện web của hệ thống đăng ký~\cite{verraregistry}. \autoref{fig:verra} minh họa danh sách dự án VCS công khai trên Verra Registry.

\begin{figure}[htbp]
    \centering
    \includegraphics[width=0.9\textwidth]{image/verra_registry.png}
    \caption[Danh sách dự án thuộc chương trình VCS trên Verra Registry]%
    {Danh sách dự án thuộc chương trình VCS trên Verra Registry~\cite{verraregistry}.}
    \label{fig:verra}
\end{figure}

Ở cấp độ quốc gia, Chương trình Phát triển của Liên Hợp Quốc (UNDP) đã cung cấp một khuôn mẫu kiến trúc đáng tham khảo thông qua National Carbon Credit Registry\index{He thong dang ky@Hệ thống đăng ký (registry)!UNDP National Carbon Registry} (phiên bản 2.0). Được phát triển dưới dạng Hàng hóa Công Kỹ thuật số (DPG)\index{DPG (Digital Public Good)} mã nguồn mở, hệ thống này hỗ trợ các quốc gia đang phát triển xây dựng hệ thống đăng ký quốc gia phục vụ Điều~6 Thỏa thuận Paris~\cite{undp2024nationalregistry}. Về mặt kỹ thuật, dữ liệu gốc của hệ thống được ghi trong hai sổ cái: Sổ cái Dự án (Project ledger) và Sổ cái Khối Tín chỉ (Credit Blocks Ledger); cơ sở dữ liệu quan hệ PostgreSQL\index{PostgreSQL} đóng vai trò cơ sở dữ liệu thứ cấp, phục vụ vận hành và truy vấn phân tích~\cite{undp2024nationalregistry}. Hệ thống được triển khai trên nền tảng đám mây AWS theo kiến trúc không máy chủ (serverless)\index{Serverless}, với các hàm AWS Lambda\index{AWS Lambda}: \texttt{national-api} xử lý các yêu cầu HTTP từ người dùng, còn \texttt{replicator} được kích hoạt bởi luồng dữ liệu Kinesis để sao chép các bản ghi từ sổ cái sang PostgreSQL~\cite{undp2024nationalregistry}.

Một điểm đáng chú ý trong thiết kế của UNDP là cách quản lý khối tín chỉ có số sê-ri liên tục (Credit Blocks)\index{Khoi tin chi@Khối tín chỉ (Credit Block)}. Khi chỉ một phần của khối tín chỉ được chuyển nhượng hoặc hủy (partial block transfer)\index{Khoi tin chi@Khối tín chỉ (Credit Block)!chuyen nhuong mot phan@chuyển nhượng một phần}, hệ thống tách dãy số sê-ri: chủ sở hữu hiện tại giữ lại phần đầu của dãy, người nhận nhận phần sau, nhờ đó vẫn truy vết được từng tín chỉ theo số sê-ri~\cite{undp2024nationalregistry}. Hệ thống cũng có chức năng tạo báo cáo theo Định dạng điện tử thỏa thuận (Agreed Electronic Format -- AEF)\index{AEF (Agreed Electronic Format)} phục vụ yêu cầu báo cáo của Điều~6.2 và hỗ trợ trao đổi dữ liệu qua tệp (ví dụ CSV)~\cite{undp2024nationalregistry}.

\subsection*{Nền tảng siêu dữ liệu phân tán: Kiến trúc CAD Trust và CDOP}
\label{subsec:cadtrust}

Các hệ thống đăng ký riêng lẻ chưa liên thông với nhau, gây khó khăn cho việc đối chiếu dữ liệu tín chỉ ở quy mô toàn cầu. Để giải quyết vấn đề này, Climate Action Data Trust (CAD Trust)\index{CAD Trust} được thành lập bởi Ngân hàng Thế giới, Chính phủ Singapore và IETA, phát triển từ nền tảng Climate Warehouse của Ngân hàng Thế giới~\cite{cadtrust2025datamodel}. CAD Trust không cấp tín chỉ và không thay thế các hệ thống đăng ký hiện có; thay vào đó, nó hoạt động như một lớp siêu dữ liệu (metadata layer)\index{Lop sieu du lieu@Lớp siêu dữ liệu (metadata layer)} phi tập trung, liên kết, tổng hợp và chuẩn hóa dữ liệu từ nhiều hệ thống đăng ký (bao gồm Verra, Gold Standard\index{Gold Standard}, Global Carbon Council, v.v.)~\cite{cadtrust2025datamodel}. Hệ sinh thái dữ liệu của CAD Trust được minh họa trong \autoref{fig:cadtrust}.

\begin{figure}[htbp]
    \centering
    \includegraphics[width=0.9\textwidth]{image/cadtrust_ecosystem.png}
    \caption[Hệ sinh thái dữ liệu của CAD Trust]%
    {Hệ sinh thái dữ liệu của CAD Trust: nguồn dữ liệu, lớp siêu dữ liệu và các ứng dụng khai thác~\cite{cadtrust2022deepdive}.}
    \label{fig:cadtrust}
\end{figure}


CAD Trust được vận hành cùng Chia Network, dựa trên chuỗi khối công khai Chia và thành phần Chia Data Layer\index{Chia Data Layer}. Thay vì tập trung dữ liệu vào một máy chủ duy nhất, mỗi hệ thống đăng ký tham gia chạy một nút Chia trên hạ tầng của mình và tự quản lý dữ liệu của mình. Mỗi lần cập nhật dữ liệu, mã băm mật mã\index{Ham bam mat ma@Hàm băm mật mã} của bản cập nhật được ghi lên chuỗi khối Chia; các tệp dữ liệu được sao chép lên lưu trữ đám mây, và một cơ sở dữ liệu MySQL được dùng để phục vụ truy vấn, phân tích~\cite{cadtrust2025datamodel}. Các bên tham gia có thể tải xuống, đồng bộ và đối chiếu dữ liệu nhận được với mã băm trên chuỗi khối để phát hiện dữ liệu bị thay đổi, mà không phải dựa vào một máy chủ trung gian duy nhất~\cite{cadtrust2025datamodel}.

Thành phần trung tâm của CAD Trust là Mô hình dữ liệu phiên bản 2.0 (Data Model Version 2.0)\index{CAD Trust!Mo hinh du lieu 2.0@Mô hình dữ liệu 2.0}, được công bố ngày 23/10/2025~\cite{cadtrust2025datamodel}. Phiên bản 1 có 9 bảng; phiên bản 2.0 giữ các bảng chung (Project, Location, Rating, Co-benefits, Estimations, Issuances, Labels, Units) và bổ sung các bảng Stakeholders, Methodologies, Programs, Validations, Verifications, Unit Labels, Project Methodologies cùng 5 bảng AEF phục vụ báo cáo theo Điều~6.2. Các hệ thống đăng ký thành viên dự kiến chuyển sang phiên bản mới trong quý I--II/2026~\cite{cadtrust2025datamodel}. Bảng~\ref{tab:cadtrust} tóm tắt các thực thể liên quan trực tiếp đến đồ án.

\begin{table}[htbp]
    \centering
    \caption{Các thực thể dữ liệu chính trong mô hình CAD Trust V2.0}
    \label{tab:cadtrust}
    \small
    \renewcommand{\arraystretch}{1.25}
    \begin{tabularx}{\textwidth}{|p{3.2cm}|X|}
        \hline
        \textbf{Thực thể dữ liệu (Bảng)} & \textbf{Thuộc tính khóa và thông tin mô tả} \\
        \hline
        Project (Dự án) & Thông tin định danh dự án (mã dự án trên hệ thống đăng ký gốc), chương trình, lĩnh vực, trạng thái; liên kết với các bảng Location (vị trí địa lý)\index{GIS} và Rating (xếp hạng)~\cite{cadtrust2025datamodel}. \\
        \hline
        Stakeholders (Các bên liên quan) & Bảng mới của V2.0, lưu các tổ chức tham gia dự án và vai trò của họ~\cite{cadtrust2025datamodel}. \\
        \hline
        Validations, Verifications & Bảng mới của V2.0, lưu thông tin các đợt validation và verification của dự án~\cite{cadtrust2025datamodel}. \\
        \hline
        Issuances (Đợt cấp) & Mỗi bản ghi là một đợt cấp tín chỉ của dự án, liên kết với Project~\cite{cadtrust2025datamodel}. \\
        \hline
        Units (Đơn vị tín chỉ) & Khối tín chỉ với dãy số sê-ri, năm thực hiện giảm phát thải (Vintage Year)\index{Vintage Year}, trạng thái (Issued, Held, Retired); với tín chỉ đã hủy, lưu thêm ngày hủy và người thụ hưởng (beneficiary)~\cite{cadtrust2025datamodel}. \\
        \hline
    \end{tabularx}
\end{table}

Mô hình dữ liệu này được thiết kế để bổ sung cho Giao thức Mở Dữ liệu Carbon (Carbon Data Open Protocol -- CDOP)\index{CDOP (Carbon Data Open Protocol)}. CDOP là một lược đồ dữ liệu mã nguồn mở do liên minh CDOP phát triển nhằm giảm tình trạng phân mảnh định dạng dữ liệu (fragmentation) và giảm chi phí hành chính. Liên minh ra mắt tháng 3/2025, do Global Carbon Market Utility, Sylvera, RMI và S\&P Global Commodity Insights đồng chủ trì, với 37 tổ chức cam kết tham gia, trong đó có Verra, Isometric và South Pole; phiên bản 1.0 được công bố ngày 23/9/2025~\cite{cdop2025press}. CDOP V1.0 tập trung vào dữ liệu trước khi cấp tín chỉ (pre-issuance) như thông tin vị trí dự án, phương pháp luận và hồ sơ công bố, còn CAD Trust tập trung vào dữ liệu sau khi cấp (post-issuance); kết hợp lại, hai mô hình bao phủ phần lớn vòng đời dữ liệu của một tín chỉ~\cite{cadtrust2025datamodel} \cite{cdop2025press}.

% ------------------------------------------------------------
\section{Khung pháp lý thị trường carbon tại Việt Nam và ứng dụng thực tiễn}
\label{sec:vietnam}

Một hệ thống quản lý tín chỉ carbon được phát triển tại Việt Nam, dù ở dạng đồ án hay sản phẩm thương mại, cần tương thích với khung pháp lý quốc gia. Văn bản pháp lý chính hiện nay là Nghị định 06/2022/NĐ-CP\index{Nghi dinh 06/2022/ND-CP@Nghị định 06/2022/NĐ-CP} quy định giảm nhẹ phát thải khí nhà kính và bảo vệ tầng ô-dôn, được sửa đổi, bổ sung bởi Nghị định 119/2025/NĐ-CP và Nghị định 83/2026/NĐ-CP~\cite{vietnam2022nghidinh06}. Từ tháng 3/2025, Bộ Tài nguyên và Môi trường được hợp nhất vào Bộ Nông nghiệp và Môi trường, nên các nhiệm vụ của Bộ trong Nghị định do Bộ Nông nghiệp và Môi trường thực hiện.

Nghị định đặt ra lộ trình phát triển thị trường carbon trong nước, bao gồm cơ chế trao đổi, bù trừ tín chỉ carbon giữa thị trường nội địa và quốc tế~\cite{vietnam2022nghidinh06}. Các cơ sở thuộc danh mục phải kiểm kê khí nhà kính được yêu cầu kiểm kê phát thải, và Bộ Nông nghiệp và Môi trường phân bổ hạn ngạch phát thải khí nhà kính\index{Han ngach phat thai@Hạn ngạch phát thải} cho các cơ sở phát thải lớn. Một hệ thống phần mềm quản lý tại Việt Nam do đó không chỉ đơn thuần theo dõi các giao dịch tự nguyện, mà còn phải cung cấp tính năng hạch toán hạn ngạch (cap-and-trade)\index{Cap-and-trade}. Thiết kế của hệ thống cần phân biệt rõ tín chỉ tự nguyện (chẳng hạn VCUs từ Verra hoặc Gold Standard) với hạn ngạch và tín chỉ thuộc thị trường trong nước do cơ quan nhà nước quản lý, đồng thời hỗ trợ lập các báo cáo định kỳ~\cite{vietnam2022nghidinh06}.

% ------------------------------------------------------------
\section{Nhiệm vụ đồ án}
\label{sec:nhiemvu}
% Khoa yêu cầu mục này nêu: tính cấp thiết, lý do chọn đề tài, ý nghĩa khoa học và thực tiễn,
% mục tiêu, đối tượng, phạm vi giới hạn. Bám sát phiếu giao nhiệm vụ (mục A1, B1–B3).

\subsection{Phát biểu bài toán}
\label{subsec:baitoan}
% Từ khảo sát ở các mục \ref{sec:dongluc}–\ref{sec:vietnam}, rút ra 3–4 vấn đề cụ thể mà hệ thống
% phải giải quyết, mỗi vấn đề trỏ về đúng mục đã khảo sát. Gợi ý:
%  - double counting (bán kép, phát hành kép, tuyên bố kép);
%  - khó kiểm chứng nguồn gốc và lịch sử sở hữu của một tín chỉ;
%  - dữ liệu phân mảnh giữa các hệ thống đăng ký, chưa theo mô hình chung (CAD Trust);
%  - việc hủy tín chỉ (retire) khó kiểm chứng công khai.
\canviet{Phát biểu bài toán}

\subsection{Mục tiêu của đồ án}
\label{subsec:muctieu}
% Lấy từ phiếu (B1–B3), viết thành các mục tiêu kiểm tra được để Chương 5 đối chiếu lại:
%  - hợp đồng thông minh: Minting (ERC-721), Trading (pool ERC-20), Retirement (burn, chống double counting);
%  - CSDL off-chain theo CAD Trust V2.0; API RESTful; giao diện Next.js + MetaMask; dashboard.
\canviet{Mục tiêu}

\subsection{Đối tượng và phạm vi}
\label{subsec:phamvi}
% Đối tượng: quy trình cấp, giao dịch, hủy tín chỉ carbon; hợp đồng thông minh trên Ethereum.
% Phạm vi: nói rõ những gì KHÔNG làm, ví dụ: chỉ triển khai trên testnet Sepolia; không xây dựng lớp
% cảm biến IoT/dMRV; không kết nối hệ thống đăng ký thật (Verra, hệ thống đăng ký quốc gia);
% dữ liệu dự án là dữ liệu mô phỏng.
\canviet{Đối tượng và phạm vi}

\subsection{Phương pháp thực hiện}
\label{subsec:phuongphap}
% Khảo sát tài liệu và hệ thống hiện có -> phân tích, thiết kế (UML) -> cài đặt -> kiểm thử trên testnet.
% Nêu công cụ chính: Solidity, Hardhat, OpenZeppelin, Node.js/Express, PostgreSQL, Next.js, ethers.js.
\canviet{Phương pháp thực hiện}

\subsection{Ý nghĩa khoa học và thực tiễn}
\label{subsec:ynghia}
\canviet{Ý nghĩa khoa học và thực tiễn}

% ------------------------------------------------------------
\section{Cấu trúc đồ án}
\label{sec:cautruc}
% Mỗi chương 2–3 câu: chương nói về gì, kết quả chính. Nên viết sau cùng, khi các chương đã xong.
\canviet{Cấu trúc đồ án}


% Chương 2 (Giai đoạn 2): blockchain, Ethereum, smart contract, chuẩn token, dApp, on/off-chain,
% nghiên cứu liên quan và khoảng trống mà đồ án giải quyết.
% ============================================================
% CHƯƠNG 2 — CƠ SỞ LÝ THUYẾT VÀ CÔNG NGHỆ  (Giai đoạn 2 của phiếu, phần tìm hiểu công nghệ)
% Nguyên tắc: chỉ trình bày kiến thức được dùng lại ở Chương 3–4. Mỗi mục nên kết thúc bằng
% 1–2 câu "áp dụng trong đồ án như thế nào".
% Hình mới của chương: figures_chap2.tex. Khóa trích dẫn gợi ý đã có sẵn trong References.bib.
% ============================================================

\chapter{CƠ SỞ LÝ THUYẾT VÀ CÔNG NGHỆ}
\label{chap:cosolythuyet}
% ============================================================
% HÌNH VẼ CHƯƠNG 2 (TikZ) — gọi bằng macro trong Chap2.tex
%   \HinhChuoiKhoi  : cấu trúc liên kết các khối          (fig:chuoikhoi)
%   \HinhLuongDApp  : luồng tương tác của dApp             (fig:luongdapp)
% ============================================================

% ---------- Cấu trúc liên kết các khối ----------
\newcommand{\HinhChuoiKhoi}{%
\begin{figure}[htbp]
\centering
\begin{tikzpicture}[
    >=Stealth,
    truong/.style={draw, minimum width=3.3cm, minimum height=0.55cm,
                   font=\footnotesize, fill=white, inner sep=1pt},
    tieude/.style={truong, fill=gray!20, font=\footnotesize\bfseries}]
  \foreach \i/\x/\ten in {0/0/N-1, 1/4.6/N, 2/9.2/N+1} {
    \node[tieude] (h\i) at (\x,0) {Khối \ten};
    \node[truong, below=0pt of h\i] (p\i) {Mã băm khối trước};
    \node[truong, below=0pt of p\i] (t\i) {Dấu thời gian};
    \node[truong, below=0pt of t\i] (m\i) {Gốc Merkle};
    \node[truong, below=0pt of m\i, minimum height=1cm, align=center] (g\i)
        {Danh sách\\giao dịch};
    \node[draw, thick, rounded corners=2pt, inner sep=3pt, fit=(h\i)(g\i)] (b\i) {};
  }
  \draw[->, thick] (p1.west) -- (b0.east |- p1.west);
  \draw[->, thick] (p2.west) -- (b1.east |- p2.west);
  \node[font=\large] at (-2.35,-1.4) {$\cdots$};
  \node[font=\large] at (11.55,-1.4) {$\cdots$};
\end{tikzpicture}
\caption{Cấu trúc liên kết giữa các khối trong chuỗi khối}
\label{fig:chuoikhoi}
\end{figure}}

% ---------- Luồng tương tác của ứng dụng phi tập trung ----------
\newcommand{\HinhLuongDApp}{%
\begin{figure}[htbp]
\centering
\begin{tikzpicture}[
    >=Stealth,
    hop/.style={draw, rounded corners=3pt, align=center, font=\footnotesize,
                text width=2.1cm, minimum height=1.1cm, fill=white},
    nhan/.style={font=\scriptsize, align=center}]
  \node[hop]                         (u)   at (0,0)     {Người dùng};
  \node[hop]                         (fe)  at (3.7,0)   {Giao diện web\\(Next.js,\\ethers.js)};
  \node[hop]                         (mm)  at (7.4,0)   {Ví\\MetaMask\\(giữ khóa riêng)};
  \node[hop, text width=2.6cm, minimum height=1.6cm, fill=gray!12]
                                     (net) at (11.3,0)  {Mạng Ethereum\\(Sepolia)\\nút RPC, EVM,\\hợp đồng thông minh};
  \node[hop]                         (be)  at (3.7,-3)  {Backend\\(Node.js, Express)};
  \node[hop]                         (db)  at (0,-3)    {CSDL off-chain\\(PostgreSQL)};

  \draw[->] (u) -- node[nhan, above] {thao tác} (fe);
  \draw[->] (fe) -- node[nhan, above] {yêu cầu ký\\(EIP-1193)} (mm);
  \draw[->] (mm) -- node[nhan, above] {giao dịch\\đã ký} (net);
  \draw[->, dashed] (fe.north) to[bend left=22]
      node[nhan, above] {đọc trạng thái on-chain (call)} (net.north);
  \draw[->] (net.south) |- node[nhan, pos=0.75, below] {sự kiện (Minted, Transfer, Retired\ldots)} (be.east);
  \draw[<->] (fe) -- node[nhan, right] {API REST} (be);
  \draw[<->] (be) -- node[nhan, above] {đọc/ghi} (db);
\end{tikzpicture}
\caption{Luồng tương tác giữa giao diện, ví MetaMask, mạng Ethereum và backend}
\label{fig:luongdapp}
\end{figure}}


% Đoạn mở đầu chương: chương trình bày nền tảng blockchain, Ethereum, hợp đồng thông minh, chuẩn token,
% dApp và cách lưu trữ on-chain/off-chain; cuối chương tổng hợp các nghiên cứu liên quan.
\canviet{Đoạn mở đầu chương}

% ------------------------------------------------------------
\section{Công nghệ chuỗi khối}
\label{sec:blockchain}

\subsection{Khái niệm và cấu trúc dữ liệu}
\label{subsec:cautrucblock}
% Ý cần có: khối, tiêu đề khối, mã băm khối trước, gốc Merkle, danh sách giao dịch;
% vì sao sửa một khối cũ làm gãy cả chuỗi (tính bất biến); sổ cái phân tán, nút.
% Trích: \cite{nakamoto2008bitcoin} \cite{zheng2017overview}. Thuật ngữ: xem Thuat_ngu.md mục 3.
% Nhớ nhắc \autoref{fig:chuoikhoi} trong đoạn văn trước khi hình xuất hiện.
\canviet{Khái niệm và cấu trúc dữ liệu chuỗi khối}

\HinhChuoiKhoi

\subsection{Cơ chế đồng thuận}
\label{subsec:dongthuan}
% Ý cần có: bằng chứng công việc (PoW) và bằng chứng cổ phần (PoS); Ethereum chuyển sang PoS
% (sự kiện The Merge, 15/9/2022). Có thể nhắc DAG (\cite{liu2024blockchain}) như một hướng khác.
% Trích: \cite{zheng2017overview} \cite{ethereum2024pos}.
\canviet{Cơ chế đồng thuận}

\subsection{Phân loại chuỗi khối}
\label{subsec:phanloai}
% Ý cần có: chuỗi khối công khai và chuỗi khối cấp phép; so sánh Ethereum với Hyperledger Fabric
% (nên làm bảng: mô hình truy cập, đồng thuận, ngôn ngữ hợp đồng, chuẩn token, chi phí, công cụ ví).
% Kết luận lý do chọn Ethereum: phiếu giao nhiệm vụ, chuẩn ERC, MetaMask, testnet công khai.
\canviet{Phân loại chuỗi khối và lý do chọn Ethereum}

% ------------------------------------------------------------
\section{Nền tảng Ethereum}
\label{sec:ethereum}

\subsection{Tài khoản và giao dịch}
\label{subsec:taikhoan}
% Ý cần có: tài khoản ngoài (EOA) và tài khoản hợp đồng; khóa riêng, địa chỉ ví; giao dịch được ký,
% nonce. Trích: \cite{buterin2014ethereum} \cite{wood2014yellow}.
\canviet{Tài khoản và giao dịch trên Ethereum}

\subsection{Máy ảo Ethereum và phí gas}
\label{subsec:evm}
% Ý cần có: EVM thực thi bytecode; gas, gas limit; phí cơ sở và phí ưu tiên theo EIP-1559.
% Đây là cơ sở cho phần đo chi phí gas ở Chương 5. Trích: \cite{wood2014yellow} \cite{eip1559}.
\canviet{Máy ảo Ethereum và phí gas}

\subsection{Mạng thử nghiệm Sepolia}
\label{subsec:sepolia}
% Ý cần có: mainnet và testnet; Sepolia, faucet; Etherscan để tra cứu giao dịch và hợp đồng.
\canviet{Mạng thử nghiệm Sepolia}

% ------------------------------------------------------------
\section{Hợp đồng thông minh}
\label{sec:smartcontract}

\subsection{Khái niệm}
\label{subsec:khainiemsc}
% Ý cần có: định nghĩa (Szabo); đặc điểm: tự thực thi, đã triển khai thì không sửa được mã;
% sự kiện (event) và nhật ký (log). Trích: \cite{szabo1997formalizing} \cite{buterin2014ethereum}.
\canviet{Khái niệm hợp đồng thông minh}

\subsection{Ngôn ngữ Solidity và thư viện OpenZeppelin}
\label{subsec:solidity}
% Ý cần có: cấu trúc một contract (biến trạng thái, hàm, modifier, event); kế thừa;
% OpenZeppelin: ERC721, ERC20, AccessControl (vai trò như VERIFIER_ROLE); Hardhat để biên dịch, kiểm thử.
% Trích: \cite{solidity2024docs} \cite{openzeppelin2024contracts}.
\canviet{Solidity và OpenZeppelin}

\subsection{Các lỗ hổng bảo mật thường gặp}
\label{subsec:baomat}
% Ý cần có: reentrancy, kiểm soát truy cập sai, tràn số (Solidity >= 0.8 đã tự kiểm tra), dùng tx.origin.
% Mỗi lỗ hổng: 1–2 câu giải thích và cách phòng trong đồ án. Công cụ quét: Slither.
% Trích: \cite{atzei2017survey}.
\canviet{Lỗ hổng bảo mật hợp đồng thông minh}

% ------------------------------------------------------------
\section{Các chuẩn token trên Ethereum}
\label{sec:token-standards}

\subsection{Chuẩn ERC-20}
\label{subsec:erc20}
% Ý cần có: token có thể thay thế; các hàm totalSupply, balanceOf, transfer, approve, transferFrom,
% allowance và sự kiện Transfer, Approval. Trích: \cite{eip20}.
\canviet{Chuẩn ERC-20}

\subsection{Chuẩn ERC-721}
\label{subsec:erc721}
% Ý cần có: token không thể thay thế (NFT); ownerOf, safeTransferFrom, tokenURI (metadata);
% sự kiện Transfer. Trích: \cite{eip721}.
\canviet{Chuẩn ERC-721}

\subsection{Lựa chọn chuẩn token cho tín chỉ carbon}
\label{subsec:chontoken}
% Ý cần có: vì sao mỗi lô tín chỉ là một NFT ERC-721 (giữ metadata riêng: dự án, vintage, số sê-ri),
% vì sao pool dùng ERC-20 (thay thế được, thanh khoản); ERC-1155 là phương án thay thế.
% Liên hệ mô hình TCO2/BCT của Toucan ở mục \ref{sec:token}. Nên có bảng so sánh ERC-20/721/1155.
\canviet{Lựa chọn chuẩn token}

% ------------------------------------------------------------
\section{Ứng dụng phi tập trung và Web3}
\label{sec:dapp}

\subsection{Kiến trúc ứng dụng phi tập trung}
\label{subsec:kientrucdapp}
% Ý cần có: các thành phần của dApp và luồng một giao dịch đi từ giao diện tới hợp đồng
% (giải thích theo \autoref{fig:luongdapp}); phần nào chạy on-chain, phần nào off-chain.
\canviet{Kiến trúc dApp}

\HinhLuongDApp

\subsection{Ví MetaMask}
\label{subsec:metamask}
% Ý cần có: ví giữ khóa riêng, khóa không rời ví; giao diện gọi ví qua chuẩn EIP-1193 (window.ethereum);
% người dùng xác nhận và ký giao dịch. Trích: \cite{eip1193}.
\canviet{Ví MetaMask}

\subsection{Thư viện ethers.js và framework Next.js}
\label{subsec:ethers}
% Ý cần có: Provider, Signer, Contract, ABI; đọc (call) và ghi (gửi giao dịch); Next.js cho giao diện.
% Trích: \cite{ethers2024docs}.
\canviet{ethers.js và Next.js}

% ------------------------------------------------------------
\section{Lưu trữ dữ liệu on-chain và off-chain}
\label{sec:onoffchain-lythuyet}
% Ý cần có: tiêu chí quyết định dữ liệu nào để on-chain (chi phí, tính công khai, kích thước, tần suất
% thay đổi) theo \cite{xu2017taxonomy}; neo mã băm (lưu keccak256 của hồ sơ trên chuỗi); IPFS định danh
% tệp theo nội dung \cite{benet2014ipfs}; đồng bộ dữ liệu bằng cách lắng nghe sự kiện.
% CAD Trust (mục \ref{subsec:cadtrust}) dùng cùng nguyên lý. Thiết kế cụ thể để ở Chương 3.
\canviet{Lưu trữ on-chain và off-chain}

% ------------------------------------------------------------
\section{Các nghiên cứu và hệ thống liên quan}
\label{sec:lienquan}
% Hai mục dưới đây chuyển nguyên từ Chương 1 (chỉ sửa thuật ngữ).
\canviet{Câu dẫn vào mục}

\subsection{Đổi mới số hóa MRV: ứng dụng blockchain và Internet vạn vật (IoT)}
\label{sec:dmrv}

Dù các hệ thống đăng ký tập trung đang dần chuẩn hóa dữ liệu, một rào cản lớn của thị trường tín chỉ carbon vẫn nằm ở khâu đo đạc, báo cáo, thẩm định (MRV) tại thực địa. MRV truyền thống phụ thuộc nhiều vào chuyên gia tư vấn có chi phí cao, kiểm toán thủ công và các báo cáo giấy theo chu kỳ dài (thường là hàng năm), gây khó khăn về nguồn lực cho các doanh nghiệp vừa và nhỏ (SMEs)\index{SMEs}~\cite{tsai2024blockchain}. Việc kết hợp Internet vạn vật (IoT)\index{IoT (Internet vạn vật)} với công nghệ sổ cái phân tán (Blockchain/DLT)\index{Blockchain}\index{DLT|see{Blockchain}} dẫn tới hướng MRV kỹ thuật số (Digital MRV)\index{MRV!Digital MRV}, trong đó dữ liệu được thu thập tự động và được dùng làm đầu vào cho hợp đồng thông minh (Smart Contract)\index{Hop dong thong minh@Hợp đồng thông minh (Smart Contract)} khi cấp tín chỉ~\cite{vaccargiu2026blockchain}.

Nghiên cứu của Vaccargiu và cộng sự (công bố tại hội thảo BoSE trong khuôn khổ IEEE SANER 2026) đề xuất một kiến trúc phần mềm hướng blockchain cho việc chứng nhận tín chỉ carbon, với ví dụ thực tế là một hệ thống quang điện (PV) công suất 100~kWp~\cite{vaccargiu2026blockchain}. Trong mô hình này, việc thu thập dữ liệu và chứng nhận được tổ chức thành 4 lớp:

\begin{enumerate}
    \item \textbf{Lớp cảm biến vật lý (Physical Sensing Layer):} Dưới cấp độ phần cứng, hệ thống triển khai 8 đồng hồ đo năng lượng thông minh 3 pha (Smart Energy Meters)\index{Dong ho do nang luong thong minh@Đồng hồ đo năng lượng thông minh}, bao phủ 24 pha điện độc lập. Các cảm biến này liên tục thu thập dữ liệu về công suất tác dụng, dòng điện, điện áp, tần số và hệ số công suất ở tần suất rất cao (1 đến 2 giây một lần)~\cite{vaccargiu2026blockchain}.
    \item \textbf{Lớp thu thập dữ liệu (Data Collection Layer):} Đóng vai trò là cổng nối, liên tục nhận luồng dữ liệu thời gian thực từ các cảm biến IoT~\cite{vaccargiu2026blockchain}.
    \item \textbf{Lớp tổng hợp dữ liệu tại biên (Data Aggregation Layer):} Gửi toàn bộ dữ liệu thô (tần suất giây) lên blockchain là không khả thi do chi phí lưu trữ lớn và tải mạng cao. Lớp tổng hợp tại biên (Edge-level aggregation)\index{Dien toan bien@Điện toán biên (Edge Computing)} tổng hợp dữ liệu theo từng phút (1.440 bản ghi mỗi ngày), sau đó gom thành các lô 5 phút để ghi lên chuỗi khối, giảm số lần ghi xuống 288 lần mỗi ngày mà vẫn giữ được tổng sản lượng điện~\cite{vaccargiu2026blockchain}.
    \item \textbf{Lớp chứng nhận và chuỗi khối (Blockchain and Certification Layer):} Dữ liệu tổng hợp được ghi vào một chuỗi khối cấp phép (permissioned blockchain)\index{Blockchain!cap phep@cấp phép (permissioned)}, cụ thể là Hyperledger Fabric\index{Hyperledger Fabric}. Tại đây, các hợp đồng thông minh kiểm tra tổng sản lượng năng lượng tái tạo theo các tiêu chí đã lập trình sẵn và thực hiện việc tạo, đăng ký tín chỉ~\cite{vaccargiu2026blockchain}. Phần giới thiệu của bài báo cho rằng hợp đồng thông minh có thể thay thế các bên trung gian như môi giới, đại lý xác minh và hệ thống đăng ký; tuy nhiên, phần thảo luận của chính bài báo vẫn giữ bước thẩm định bởi bên thứ ba được công nhận và nhấn mạnh rằng blockchain là công cụ hỗ trợ, không thay thế quá trình chứng nhận~\cite{vaccargiu2026blockchain}.
\end{enumerate}

\HinhVaccargiu

\autoref{fig:vaccargiu} minh họa tổng quan kiến trúc 4 lớp nêu trên.

Theo nhóm tác giả, cách kết hợp này giúp hồ sơ giảm phát thải khó bị sửa đổi, minh bạch và có thể được bên thứ ba kiểm tra. Ngoài giao dịch tín chỉ, dữ liệu này còn có thể phục vụ báo cáo bền vững của doanh nghiệp theo Chỉ thị Báo cáo Bền vững Doanh nghiệp của Liên minh Châu Âu (Directive EU 2022/2464)\index{CSRD (Directive EU 2022/2464)}, góp phần hạn chế việc tuyên bố sai về môi trường (greenwashing)\index{Greenwashing}~\cite{vaccargiu2026blockchain}. Bài báo nêu nguyên tắc một tín chỉ chỉ có thể được bán hoặc được dùng cho tuyên bố nội bộ, không thể cả hai, và cho rằng sổ cái phân tán khiến việc tuyên bố cùng một lượng giảm phát thải nhiều lần trở nên khó khăn hơn; tuy nhiên, bài báo chưa mô tả cụ thể cơ chế khóa hay hủy tín chỉ trong hợp đồng thông minh~\cite{vaccargiu2026blockchain}. Đây là một khác biệt với đồ án, vốn thiết kế cụ thể bước VVB phê duyệt trước khi mint và hàm hủy tín chỉ (mục \ref{subsec:khoangtrong}).

Ngoài ra, nghiên cứu của Liu và cộng sự (2024) trên tạp chí \textit{Energy} (Elsevier) đề xuất dùng chuỗi khối dạng Đồ thị Có hướng Không chu trình (Directed Acyclic Graph -- DAG)\index{DAG (Directed Acyclic Graph)} cho giao dịch chứng chỉ năng lượng tái tạo (RECs)\index{RECs} giữa nhà cung cấp và người tiêu dùng. Khác với blockchain dùng Proof-of-Work\index{Co che dong thuan@Cơ chế đồng thuận!Proof-of-Work} (tiêu thụ nhiều điện năng) hoặc Proof-of-Stake\index{Co che dong thuan@Cơ chế đồng thuận!Proof-of-Stake}, trong cấu trúc DAG các giao dịch không cần được gom vào khối tuần tự; mỗi giao dịch mới xác nhận một số giao dịch trước đó. Theo kết quả mô phỏng của nhóm tác giả, so với cơ chế Proof-of-Stake, phương án đề xuất giảm 41\% thời gian giao dịch và 65\% năng lượng tiêu thụ~\cite{liu2024blockchain}. Để bảo vệ thông tin của người tiêu dùng, hệ thống sử dụng định danh phi tập trung (DID), chữ ký số, mật khẩu dùng một lần (OTP) và các tài khoản trung gian để che giấu mẫu hình giao dịch~\cite{liu2024blockchain}. Hướng tiếp cận này phù hợp với các bên tham gia không muốn công khai thông tin giao dịch của mình trên sổ cái~\cite{liu2024blockchain}.

\subsection{Token hóa (tokenization) và giao thức tài chính phi tập trung (DeFi)}
\label{sec:token}

Trong khi các giải pháp như Hyperledger Fabric hướng tới môi trường doanh nghiệp cấp phép (permissioned), các dự án Web3\index{Web3} và Tổ chức Tự trị Phi tập trung (DAO)\index{DAO} tiếp cận thị trường carbon theo hướng token hóa tài sản thực (Real World Assets -- RWA Tokenization)\index{Token hoa@Token hóa (Tokenization)}\index{RWA|see{Token hóa}}. Token hóa tạo ra các token trên chuỗi khối đại diện cho tín chỉ carbon đã được cấp trên hệ thống đăng ký truyền thống, kèm theo các thuộc tính của tín chỉ đó~\cite{toucan2025docs}.

Toucan Protocol\index{Toucan Protocol} và KlimaDAO\index{KlimaDAO} là hai trong những nền tảng đầu tiên xây dựng cầu nối giữa tín chỉ carbon và tài chính phi tập trung (DeFi)\index{DeFi}, đưa các tín chỉ đã được cấp (như VCUs từ Verra) lên các chuỗi khối tương thích EVM\index{EVM} như Polygon và Celo~\cite{toucan2025docs}. Toucan Protocol dùng nhiều loại token khác nhau để vừa giữ được thông tin riêng của từng dự án, vừa tạo ra token có thể thay thế cho nhau để tăng thanh khoản~\cite{toucan2025docs}. Quy trình token hóa và hủy tín chỉ qua Toucan gồm các bước sau:

\begin{enumerate}
    \item \textbf{Cầu nối từ Verra và Batch NFT (ERC-721):}\index{Toucan Protocol!Batch NFT} Người dùng tạo một Batch NFT (ERC-721)\index{NFT} trên Polygon, sau đó hủy lô tín chỉ trên Verra Registry và ghi mã định danh của Batch NFT vào ghi chú hủy. Mỗi lô chỉ được gồm tín chỉ của cùng một dự án và cùng năm vintage. Số sê-ri của Verra được ghi vào NFT và được người xác minh của Toucan phê duyệt~\cite{toucan2025docs}.
    \item \textbf{Tạo TCO2 (ERC-20):}\index{Toucan Protocol!TCO2} Batch NFT đã được duyệt được chia nhỏ thành token TCO2. TCO2 là token ERC-20\index{ERC-20}, mỗi đợt cấp (vintage) của một dự án có một hợp đồng TCO2 riêng giữ siêu dữ liệu của dự án; 1 TCO2 tương ứng 1 tín chỉ (1 tCO$_2$e)~\cite{toucan2025docs}. Cầu nối từ Verra hiện đã ngừng hoạt động vì Verra không còn chấp nhận việc token hóa tín chỉ theo cách này~\cite{toucan2025docs}.
    \item \textbf{Nạp vào pool và đổi lại (Pool Tokens -- ERC-20):} Việc giao dịch trực tiếp hàng nghìn loại TCO2 khác nhau làm thanh khoản bị phân tán. Toucan dùng các hợp đồng pool: các TCO2 đáp ứng tiêu chí chấp nhận của pool (ví dụ năm vintage, loại dự án) có thể được nạp vào pool để nhận token ERC-20 có thể thay thế như Base Carbon Tonne (BCT)\index{Toucan Protocol!BCT (Base Carbon Tonne)} hoặc Nature Carbon Tonne (NCT)\index{Toucan Protocol!NCT (Nature Carbon Tonne)}, mỗi token tương ứng 1 tấn CO$_2$e và có thể giao dịch trên các sàn phi tập trung (DEX)\index{DEX}. Ngược lại, người giữ pool token có thể đổi lại (redeem) lấy TCO2 của một dự án cụ thể~\cite{toucan2025docs}.
    \item \textbf{Hủy tín chỉ (retire):} Người giữ TCO2 gọi hàm \texttt{retire()} của hợp đồng TCO2 để hủy tín chỉ; lượng token tương ứng bị đốt (burn) và loại khỏi lưu thông. Hàm \texttt{retireAndMintCertificate()} vừa hủy tín chỉ vừa tạo NFT chứng nhận hủy (RetirementCertificate)\index{Huy tin chi@Hủy tín chỉ (retirement)!tren chuoi@trên chuỗi khối}, ghi tên bên hủy, người thụ hưởng và thông điệp hủy; các thông tin này được ghi công khai trên chuỗi khối~\cite{toucan2025docs}. Toucan cũng cung cấp hợp đồng \texttt{OffsetHelper}\index{Toucan Protocol!OffsetHelper} để gộp các bước hoán đổi, đổi lại và hủy trong một giao dịch~\cite{toucan2025docs}.
\end{enumerate}
% [GÓP Ý] Đã bỏ các chi tiết lấy từ bài blog Medium nhưng không đối chiếu được với tài liệu chính thức
% của Toucan: tên hàm retireSpecificProject(), phí đổi 10% của NCT, chức năng đa chuỗi qua Hyperlane.

\HinhToucan

\autoref{fig:toucan} tóm tắt luồng token hóa và hủy tín chỉ theo giao thức Toucan.

Tại Đài Loan, nghiên cứu của Tsai (2024) xây dựng một nền tảng giao dịch tín chỉ carbon trên Ethereum\index{Ethereum} cho các doanh nghiệp vừa và nhỏ, nhằm hỗ trợ đáp ứng Cơ chế điều chỉnh biên giới carbon (CBAM)\index{CBAM} của EU -- có giai đoạn áp dụng chính thức từ năm 2026 -- và dự luật Clean Competition Act của Hoa Kỳ~\cite{tsai2024blockchain}. Hệ thống dùng giao diện React.js\index{React.js} kết nối ví MetaMask\index{MetaMask}, còn các chức năng chính như phát hành và hủy tín chỉ được thực hiện bằng hợp đồng thông minh Solidity\index{Solidity}. Để giảm chi phí gas, hồ sơ doanh nghiệp có dung lượng lớn được lưu ngoài chuỗi (off-chain)\index{Off-chain}, chỉ mã băm và thông tin chuyển quyền sở hữu được ghi lên chuỗi (on-chain)\index{On-chain}, tạo thành một mô hình lai cho đối tượng SMEs~\cite{tsai2024blockchain}. Trong một thí nghiệm có đối chứng với hệ thống tập trung, tác giả ghi nhận nền tảng giúp giảm thời gian và chi phí giao dịch, đồng thời dùng mô hình Kano để đánh giá mức độ hài lòng của người dùng~\cite{tsai2024blockchain}.

\subsection{Đánh giá và khoảng trống mà đồ án giải quyết}
\label{subsec:khoangtrong}
% Nên có bảng so sánh: Verra Registry, UNDP National Carbon Registry, CAD Trust, Toucan Protocol,
% Vaccargiu và cộng sự, nền tảng cho SMEs, và hệ thống của đồ án; tiêu chí gợi ý: nền tảng
% (tập trung / blockchain nào), quyền sở hữu ghi on-chain, cơ chế chống double counting,
% mô hình dữ liệu theo chuẩn, mã nguồn mở. Sau bảng, nêu rõ điểm đồ án kế thừa và điểm khác.
\canviet{Đánh giá và khoảng trống}

% ------------------------------------------------------------
\section{Kết luận chương}
% 3–5 câu: các kiến thức chính và lựa chọn công nghệ được mang sang Chương 3.
\canviet{Kết luận chương 2}


% Chương 3 (Giai đoạn 2 → 4): yêu cầu, ca sử dụng, kiến trúc, phân tách on/off-chain, vòng đời tín chỉ;
% thiết kế contract, CSDL, API, biểu đồ tuần tự bổ sung ở Giai đoạn 3–4.
% ============================================================
% CHƯƠNG 3 — PHÂN TÍCH VÀ THIẾT KẾ HỆ THỐNG  (Giai đoạn 2 của phiếu, phần thiết kế kiến trúc)
% Các mục thiết kế chi tiết (hợp đồng thông minh, CSDL, API, biểu đồ tuần tự) thuộc Giai đoạn 3–4,
% bổ sung vào cuối chương khi làm tới.
% Hình mới của chương: figures_chap3.tex.
% ============================================================

\chapter{PHÂN TÍCH VÀ THIẾT KẾ HỆ THỐNG}
\label{chap:thietke}
% ============================================================
% HÌNH VẼ CHƯƠNG 3 (TikZ) — gọi bằng macro trong Chap3.tex
%   \HinhUseCase     : biểu đồ ca sử dụng tổng quát        (fig:usecase)
%   \HinhOnOffChain  : phân tách dữ liệu on/off-chain       (fig:onoffchain)
%   \HinhTrangThai   : biểu đồ trạng thái của lô tín chỉ    (fig:trangthai)
% Đây là bản đề xuất theo phiếu giao nhiệm vụ: sửa lại cho khớp thiết kế thực tế của bạn.
% ============================================================

% Hình người (tác nhân UML)
\tikzset{
  tacnhan/.pic={
    \draw[thick] (0,0.62) circle (0.15);
    \draw[thick] (0,0.47) -- (0,0.05);
    \draw[thick] (-0.25,0.33) -- (0.25,0.33);
    \draw[thick] (0,0.05) -- (-0.18,-0.28);
    \draw[thick] (0,0.05) -- (0.18,-0.28);
  }
}

% ---------- Biểu đồ ca sử dụng tổng quát ----------
\newcommand{\HinhUseCase}{%
\begin{figure}[htbp]
\centering
\begin{tikzpicture}[font=\footnotesize]
  % Khung hệ thống
  \draw[thick] (3.0,-4.3) rectangle (9.8,4.6);
  \node[font=\footnotesize\bfseries] at (6.4,4.25) {Hệ thống quản lý tín chỉ carbon};
  % Ca sử dụng: tên nút / tung độ / nội dung
  \foreach \ten/\y/\noidung in {
      uc1/3.5/{Thẩm định, phê duyệt và mint tín chỉ},
      uc2/2.6/{Đăng ký dự án},
      uc3/1.7/{Gửi yêu cầu cấp tín chỉ},
      uc4/0.8/{Nạp tín chỉ vào pool, đổi token},
      uc5/-0.1/{Mua bán, chuyển nhượng tín chỉ},
      uc6/-1.0/{Hủy tín chỉ để bù trừ (retire)},
      uc7/-1.9/{Tra cứu tín chỉ, chứng nhận hủy},
      uc8/-2.8/{Theo dõi dashboard hạn ngạch},
      uc9/-3.7/{Quản lý tài khoản, phân quyền}} {
    \draw (6.4,\y) ellipse (3.2cm and 0.38cm);
    \node at (6.4,\y) {\noidung};
    \coordinate (\ten-w) at (3.2,\y);
    \coordinate (\ten-e) at (9.6,\y);
  }
  % Tác nhân bên trái (tên ở dưới)
  \pic at (1.2,3.1) {tacnhan};   \node[align=center] at (1.2,2.55) {VVB};
  \pic at (1.2,0.8) {tacnhan};   \node[align=center] at (1.2,0.25) {Chủ dự án};
  % Tác nhân bên phải (tên ở bên phải để không đè lên đường nối)
  \pic at (11.2,0.5) {tacnhan};  \node[align=left, anchor=west] at (11.5,0.55) {Doanh nghiệp\\mua};
  \pic at (11.2,-3.0) {tacnhan}; \node[align=left, anchor=west] at (11.5,-2.95) {Cơ quan\\quản lý};
  % Liên kết
  \draw (1.5,3.4) -- (uc1-w);
  \foreach \u in {uc2,uc3,uc4,uc5} \draw (1.5,1.1) -- (\u-w);
  \foreach \u in {uc4,uc5,uc6,uc7} \draw (10.9,0.8) -- (\u-e);
  \foreach \u in {uc7,uc8,uc9} \draw (10.9,-2.7) -- (\u-e);
\end{tikzpicture}
\caption{Biểu đồ ca sử dụng tổng quát của hệ thống}
\label{fig:usecase}
\end{figure}}

% ---------- Phân tách dữ liệu on-chain / off-chain ----------
\newcommand{\HinhOnOffChain}{%
\begin{figure}[htbp]
\centering
\begin{tikzpicture}[
    >=Stealth,
    o/.style={draw, rounded corners=2pt, align=left, font=\footnotesize,
              text width=4.7cm, inner sep=4pt, fill=white},
    nhan/.style={font=\scriptsize, align=center, text width=2.6cm}]
  % Off-chain
  \node[o] (files) at (2.6,1.6) {\textbf{Kho tệp}\\
      IPFS hoặc lưu trữ đối tượng:\\
      PDD, báo cáo thẩm định,\\ tệp metadata JSON};
  \node[o, below=5mm of files] (db) {\textbf{CSDL quan hệ}\\
      PostgreSQL, theo mô hình\\ CAD Trust V2.0:\\
      Project, Stakeholders,\\ Verifications, Issuances, Units;\\
      tài khoản ứng dụng (bảng riêng);\\
      lịch sử giao dịch đồng bộ từ sự kiện};
  \node[draw, thick, dashed, rounded corners=4pt, inner sep=6pt, fit=(files)(db),
        label={[font=\footnotesize\bfseries]above:Off-chain (backend)}] (off) {};
  % On-chain
  \node[o, fill=gray!10] (nft) at (11.0,1.6) {\textbf{Hợp đồng NFT tín chỉ (ERC-721)}\\
      tokenId, chủ sở hữu, số sê-ri,\\ mã băm hồ sơ, URI metadata};
  \node[o, fill=gray!10, below=4mm of nft] (pool) {\textbf{Hợp đồng pool (ERC-20)}\\
      số dư token,\\ NFT đang khóa trong pool};
  \node[o, fill=gray!10, below=4mm of pool] (ret) {\textbf{Module Retirement}\\
      sự kiện Retired, người thụ hưởng};
  \node[draw, thick, rounded corners=4pt, inner sep=6pt, fit=(nft)(pool)(ret),
        label={[font=\footnotesize\bfseries]above:On-chain (Ethereum Sepolia)}] (on) {};
  % Luồng dữ liệu
  \draw[->, thick] (files.east) -- node[nhan, above] {mã băm hồ sơ (keccak256) và URI khi mint}
      (on.west |- files.east);
  \draw[->, thick] (on.west |- db.east) -- node[nhan, below] {sự kiện Minted, Transfer, Retired: backend đồng bộ}
      (db.east);
\end{tikzpicture}
\caption{Phân tách dữ liệu giữa phần on-chain và off-chain của hệ thống}
\label{fig:onoffchain}
\end{figure}}

% ---------- Biểu đồ trạng thái của một lô tín chỉ ----------
\newcommand{\HinhTrangThai}{%
\begin{figure}[htbp]
\centering
\begin{tikzpicture}[
    >=Stealth,
    tt/.style={draw, rounded corners=6pt, align=center, font=\footnotesize,
               minimum height=0.9cm, inner xsep=5pt, fill=white},
    nhan/.style={font=\scriptsize, align=center}]
  \node[circle, fill=black, inner sep=3pt] (bd) at (0,0) {};
  \node[tt] (nhap)   at (1.4,0)   {Nháp};
  \node[tt] (cho)    at (4.5,0)   {Chờ thẩm định};
  \node[tt] (tuchoi) at (4.5,-2.3) {Bị từ chối};
  \node[tt] (dacap)  at (9.0,0)   {Đã cấp\\(NFT ERC-721)};
  \node[tt] (pool)   at (9.0,-2.6) {Trong pool\\(ERC-20)};
  \node[tt] (dahuy)  at (12.6,0)   {Đã hủy\\(retired)};
  \node[circle, draw, inner sep=4pt] (kt) at (14.1,0) {};
  \node[circle, fill=black, inner sep=2pt] at (kt) {};

  \draw[->] (bd) -- (nhap);
  \draw[->] (nhap) -- node[nhan, above] {gửi\\hồ sơ} (cho);
  \draw[->] (cho) -- node[nhan, above] {VVB duyệt,\\mint NFT} (dacap);
  \draw[->] (cho) -- node[nhan, right] {VVB từ chối} (tuchoi);
  \draw[->] (tuchoi.west) -| node[nhan, pos=0.3, below] {bổ sung hồ sơ} (nhap.south);
  \draw[->] (dacap) -- node[nhan, above] {retire\\(burn)} (dahuy);
  \draw[->] (dahuy) -- (kt);
  \draw[->] ([xshift=-5mm]dacap.south) -- node[nhan, left] {nạp vào pool} ([xshift=-5mm]pool.north);
  \draw[->] ([xshift=5mm]pool.north) -- node[nhan, right] {đổi lại\\(redeem)} ([xshift=5mm]dacap.south);
  \draw[->] ([xshift=-4mm]dacap.north) .. controls +(-0.3,0.9) and +(0.3,0.9) ..
      node[nhan, above] {chuyển nhượng} ([xshift=4mm]dacap.north);
\end{tikzpicture}
\caption{Biểu đồ trạng thái của một lô tín chỉ carbon trong hệ thống}
\label{fig:trangthai}
\end{figure}}


% Đoạn mở đầu chương: từ bài toán ở mục \ref{subsec:baitoan} và công nghệ ở Chương 2,
% chương này xác định yêu cầu, rồi thiết kế kiến trúc và luồng dữ liệu on-chain/off-chain.
\canviet{Đoạn mở đầu chương}

% ------------------------------------------------------------
\section{Phân tích yêu cầu}
\label{sec:yeucau}

\subsection{Các tác nhân}
\label{subsec:tacnhan}
% Nên làm bảng: Tác nhân | Vai trò nghiệp vụ | Vai trò trong hợp đồng (role).
% Gợi ý: Chủ dự án; VVB (VERIFIER_ROLE); Doanh nghiệp mua (người hủy tín chỉ);
% Cơ quan quản lý (DEFAULT_ADMIN_ROLE). Người xem công khai chỉ tra cứu.
\canviet{Các tác nhân}

\subsection{Biểu đồ ca sử dụng tổng quát}
\label{subsec:usecase}
% Giải thích \autoref{fig:usecase}. Sau đó đặc tả 3 ca sử dụng chính bằng bảng
% (Tên, Tác nhân, Tiền điều kiện, Luồng chính, Luồng thay thế, Hậu điều kiện):
% Gửi yêu cầu cấp tín chỉ; Thẩm định, phê duyệt và mint; Hủy tín chỉ.
% Mọi ca sử dụng ghi on-chain đều cần ``Kết nối ví MetaMask'' làm tiền điều kiện.
\canviet{Mô tả biểu đồ ca sử dụng}

\HinhUseCase

\subsection{Yêu cầu chức năng}
\label{subsec:ycchucnang}
% Bảng: Mã (YC-01, YC-02, ...) | Tên yêu cầu | Tác nhân | Mô tả ngắn. Mã này sẽ dùng lại ở Chương 5
% để đối chiếu kết quả kiểm thử.
\canviet{Yêu cầu chức năng}

\subsection{Yêu cầu phi chức năng}
\label{subsec:ycphichucnang}
% Gợi ý nhóm: toàn vẹn (mỗi số sê-ri chỉ mint một lần; tín chỉ đã hủy không chuyển được);
% bảo mật (phân quyền theo role, chống reentrancy); minh bạch (tra cứu công khai);
% chi phí gas; tương thích dữ liệu với CAD Trust V2.0; hiệu năng API.
\canviet{Yêu cầu phi chức năng}

\section{Đề xuất kiến trúc tổng thể cho hệ thống quản lý tín chỉ carbon}
\label{sec:kientruc}

Từ kết quả khảo sát ở Chương~\ref{chap:khaosat} (Điều~6.2/6.4 Thỏa thuận Paris, Nghị định 06/2022, các hệ thống đăng ký UNDP và Verra, nền tảng siêu dữ liệu CAD Trust) và các nghiên cứu liên quan ở mục \ref{sec:lienquan} (MRV số hóa, Toucan Protocol), đồ án đề xuất một kiến trúc lai (Hybrid Architecture)\index{Kien truc lai@Kiến trúc lai (Hybrid Architecture)} kết hợp phần off-chain và on-chain. Dữ liệu có dung lượng lớn hoặc thường thay đổi (hồ sơ dự án, siêu dữ liệu, tài khoản người dùng) được lưu off-chain để giảm chi phí và hỗ trợ truy vấn nhanh; các thao tác cần kiểm chứng công khai (cấp, chuyển nhượng và hủy tín chỉ) được thực hiện on-chain. Hệ thống được tổ chức theo kiến trúc phân lớp\index{Kien truc phan lop@Kiến trúc phân lớp (Layered Architecture)} gồm 4 lớp như sau.

\subsection*{Lớp 1: Lớp nguồn dữ liệu và MRV (Data Source \& MRV Layer)}
\label{subsec:lop1}

Lớp này cung cấp dữ liệu đầu vào cho quá trình cấp tín chỉ: hồ sơ thiết kế dự án (PDD)\index{PDD (Project Design Document)}, báo cáo giám sát và báo cáo thẩm định.
\begin{itemize}
    \item Trong phạm vi đồ án, Chủ dự án nhập thông tin và tải hồ sơ lên qua giao diện web; VVB tải lên báo cáo thẩm định; số liệu giám sát là dữ liệu mô phỏng.
    \item Việc thu thập dữ liệu tự động bằng cảm biến IoT và tổng hợp tại biên theo mô hình của Vaccargiu và cộng sự~\cite{vaccargiu2026blockchain} không được cài đặt trong đồ án và được đưa vào phần hướng phát triển. Kiến trúc để sẵn điểm tích hợp: dữ liệu tổng hợp từ thiết bị có thể được gửi vào Lớp~2 qua API REST.
\end{itemize}

\subsection*{Lớp 2: Lớp dịch vụ và lưu trữ off-chain (Off-chain Service \& Storage Layer)}
\label{subsec:lop2}

\begin{itemize}
    \item \textbf{Backend:} Node.js và Express, cung cấp API REST\index{REST API}, xác thực người dùng và kiểm tra tính hợp lệ của hồ sơ trước khi chuyển cho VVB.
    \item \textbf{Cơ sở dữ liệu quan hệ:} PostgreSQL\index{PostgreSQL}, thiết kế dựa trên các thực thể của mô hình CAD Trust V2.0 (Bảng~\ref{tab:cadtrust}): Project, Stakeholders, Validations/Verifications, Issuances và Units~\cite{cadtrust2025datamodel}. Đồ án dùng thống nhất tên Unit cho khối tín chỉ. Mỗi Unit tương ứng một NFT có dãy số sê-ri cố định và không được tách khi chuyển nhượng; nhu cầu giao dịch từng phần được đáp ứng qua pool ERC-20 ở Lớp~3, thay cho cách tách dãy số sê-ri của UNDP~\cite{undp2024nationalregistry}. Tài khoản đăng nhập và phân quyền của ứng dụng được lưu ở bảng riêng, tách khỏi các bảng theo mô hình CAD Trust.
    \item \textbf{Lưu trữ tệp:} tệp PDD, báo cáo thẩm định và tệp metadata JSON được lưu trên IPFS\index{IPFS} hoặc dịch vụ lưu trữ đối tượng (như AWS S3\index{AWS S3}); cơ sở dữ liệu chỉ lưu đường dẫn và mã băm của tệp.
    \item \textbf{Đồng bộ sự kiện:} một tiến trình lắng nghe sự kiện (event listener) đọc các sự kiện \texttt{Minted}, \texttt{Transfer}, \texttt{Retired} từ hợp đồng để cập nhật cơ sở dữ liệu. Quyền sở hữu và trạng thái hủy lấy dữ liệu on-chain làm gốc; cơ sở dữ liệu off-chain là bản sao phục vụ truy vấn, tương tự vai trò của PostgreSQL trong hệ thống của UNDP~\cite{undp2024nationalregistry}.
\end{itemize}

\subsection*{Lớp 3: Lớp hợp đồng thông minh (On-chain Smart Contract Layer)}
\label{subsec:lop3}

Các hợp đồng được viết bằng Solidity, sử dụng thư viện OpenZeppelin và triển khai trên mạng thử nghiệm Ethereum Sepolia theo phiếu giao nhiệm vụ. Việc so sánh với chuỗi khối cấp phép như Hyperledger Fabric được trình bày ở mục \ref{subsec:phanloai}.
\begin{itemize}
    \item \textbf{Module Minting (Phát hành):} VVB (vai trò \texttt{VERIFIER\_ROLE}) thẩm định hồ sơ ở ngoài chuỗi; khi phê duyệt, VVB gọi hàm mint từ ví của mình. Hợp đồng kiểm tra quyền của bên gọi và kiểm tra dãy số sê-ri chưa từng được cấp, sau đó tạo một NFT ERC-721\index{ERC-721} cho lô tín chỉ, lưu số lượng tCO$_2$e, dãy số sê-ri, mã băm hồ sơ và URI metadata. Hợp đồng không tự đánh giá nội dung hồ sơ; độ chính xác của số liệu dựa vào kết quả thẩm định của VVB. Cách tách giữa NFT đại diện cho lô và token có thể thay thế tham khảo mô hình Batch NFT và TCO2 của Toucan~\cite{toucan2025docs}.
    \item \textbf{Module Pool (Giao dịch):} Chủ sở hữu có thể nạp NFT vào hợp đồng pool; NFT bị khóa trong pool và pool đúc lượng token ERC-20\index{ERC-20} tương ứng (1 token = 1 tCO$_2$e) để giao dịch từng phần. Thao tác đổi lại (redeem) đốt lượng token tương ứng để rút một lô ra khỏi pool~\cite{toucan2025docs}.
    \item \textbf{Module Retirement (Hủy tín chỉ):} Khi người dùng chọn ``Bù trừ phát thải'', hàm \texttt{retire()} đốt (burn) token và phát sự kiện \texttt{Retired} ghi số lượng, dãy số sê-ri và người thụ hưởng. Vì token không còn tồn tại, tín chỉ đã hủy không thể bị chuyển nhượng hoặc hủy lại trong hệ thống. Cơ chế này không ngăn được việc cùng một lượng giảm phát thải được đăng ký ở một hệ thống khác (phát hành kép giữa các hệ thống); vấn đề đó cần đến liên thông dữ liệu như CAD Trust. Cách đốt token (hàm \texttt{\_burn} của OpenZeppelin hay chuyển tới địa chỉ đốt) và cách ghi nhận phần tín chỉ hủy từ pool được chọn ở mục thiết kế hợp đồng thông minh.
\end{itemize}

\subsection*{Lớp 4: Lớp giao diện ứng dụng và API (Presentation \& API Layer)}
\label{subsec:lop4}

\begin{itemize}
    \item \textbf{Giao diện người dùng (Frontend):} Xây dựng bằng Next.js và ethers.js, kết nối ví MetaMask để người dùng ký các giao dịch mint, nạp pool, chuyển nhượng và hủy tín chỉ; khóa riêng không rời khỏi ví. Cách kết hợp giao diện web với MetaMask tương tự nền tảng của Tsai~\cite{tsai2024blockchain}.
    \item \textbf{Dashboard và tra cứu:} Hiển thị tín chỉ đang nắm giữ, đã hủy và lịch sử giao dịch. Với doanh nghiệp, dashboard so sánh lượng tín chỉ đã hủy với lượng phát thải hoặc hạn ngạch đã khai báo (dữ liệu mô phỏng). Người xem công khai có thể tra cứu tín chỉ theo số sê-ri, kèm liên kết tới giao dịch trên Etherscan.
    \item \textbf{API liên thông:} API REST xuất dữ liệu dự án và tín chỉ theo cấu trúc trường của CAD Trust V2.0 (Project, Issuances, Units) để có thể trao đổi với hệ thống bên ngoài~\cite{cadtrust2025datamodel}. Việc xuất báo cáo theo Định dạng Điện tử Thỏa thuận (AEF) của Điều~6.2 như hệ thống của UNDP~\cite{undp2024nationalregistry} là hướng mở rộng, không thuộc phạm vi cài đặt.
\end{itemize}

\HinhKienTruc

\autoref{fig:kientruc} trình bày kiến trúc tổng thể hệ thống đề xuất.

Kiến trúc trên kết hợp một cơ sở dữ liệu off-chain có cấu trúc dựa trên mô hình CAD Trust V2.0 với các hợp đồng thông minh ERC-721/ERC-20 cho ba thao tác cấp, giao dịch và hủy tín chỉ. Phần on-chain giúp quyền sở hữu và trạng thái hủy có thể được kiểm chứng công khai; phần off-chain giữ chi phí lưu trữ thấp và hỗ trợ truy vấn. Hệ thống được xây dựng ở mức nguyên mẫu trên mạng thử nghiệm và chưa kết nối với hệ thống đăng ký chính thức nào.

% ------------------------------------------------------------
\section{Phân tách dữ liệu on-chain và off-chain}
\label{sec:onoffchain}
% Giải thích \autoref{fig:onoffchain}. Nên kèm bảng: Dữ liệu | Nơi lưu | Lý do (chi phí, công khai,
% kích thước, cần bất biến hay không). Dựa trên tiêu chí ở mục \ref{sec:onoffchain-lythuyet}.
% Nêu cơ chế giữ hai phía nhất quán: mã băm hồ sơ lưu trong NFT; backend lắng nghe sự kiện để cập nhật CSDL.
\canviet{Phân tách dữ liệu on-chain và off-chain}

\HinhOnOffChain

% ------------------------------------------------------------
\section{Vòng đời tín chỉ carbon trong hệ thống}
\label{sec:vongdoi}
% Giải thích \autoref{fig:trangthai}: mỗi chuyển trạng thái ứng với hàm nào của hợp đồng / API nào.
% Chỉ ra ba điểm chặn double counting: (1) mint kiểm tra số sê-ri chưa tồn tại (phát hành kép);
% (2) NFT nạp vào pool bị khóa, không bán riêng được (bán kép); (3) retire burn token và ghi người
% thụ hưởng (tuyên bố kép).
\canviet{Vòng đời tín chỉ}

\HinhTrangThai

% ------------------------------------------------------------
% CÁC MỤC BỔ SUNG Ở GIAI ĐOẠN 3–4:
%   - Thiết kế hợp đồng thông minh: biểu đồ lớp các contract, hàm, event, modifier, role.
%   - Thiết kế CSDL off-chain: ERD theo CAD Trust V2.0 (Project, Stakeholders, Validations/Verifications, Issuances, Units)
%     và bảng tài khoản ứng dụng.
%   - Thiết kế API: bảng endpoint.
%   - Biểu đồ tuần tự: cấp tín chỉ (mint), giao dịch qua pool, hủy tín chỉ (retire).
% ------------------------------------------------------------

\section{Kết luận chương}
% 3–5 câu: yêu cầu đã xác định, kiến trúc đã chọn, những gì Chương 4 sẽ cài đặt.
\canviet{Kết luận chương 3}


% Chương 4: môi trường và công nghệ, cài đặt smart contract, backend, frontend,
% triển khai lên testnet.
\chapter{XÂY DỰNG HỆ THỐNG}
Nội dung chương 4...

% Chương 5: kịch bản và kết quả kiểm thử, đo chi phí gas, kịch bản chống double counting,
% đối chiếu kết quả với yêu cầu ở chương 3.
\chapter{KẾT QUẢ THỰC NGHIỆM}
Nội dung chương 5...

% ================= KẾT LUẬN =================
% Theo mẫu: không đưa phương trình, biểu đồ, bảng biểu vào phần kết luận; trình bày
% trung thực mức độ đáp ứng yêu cầu đề bài, hạn chế, kiến thức và kỹ năng đã đạt được.
\chapter*{Kết luận}
\phantomsection
\addcontentsline{toc}{chapter}{KẾT LUẬN}

\section*{Kết luận}
Nội dung kết luận...

\section*{Hướng phát triển của đồ án trong tương lai}
Nội dung hướng phát triển...

% ================= PHỤ LỤC (đánh theo A, B, C...) =================
% Phụ lục không được dày hơn phần chính của đồ án. Với mã nguồn, chỉ đưa các đoạn mã
% quan trọng và đường dẫn GitHub, địa chỉ contract trên Etherscan.
\appendix
\titleformat{\chapter}[block]
    {\centering\normalsize\bfseries}
    {PHỤ LỤC \thechapter.}{0.5em}{\MakeUppercase}
\addtocontents{toc}{\protect\renewcommand{\protect\cftchappresnum}{PHỤ LỤC }}
\addtocontents{toc}{\protect\setlength{\protect\cftchapnumwidth}{3.1cm}}

\chapter{MÃ NGUỒN CHƯƠNG TRÌNH}
% Ví dụ chèn mã: \lstinputlisting[language=Solidity,caption={...}]{contracts/CarbonCredit.sol}

\chapter{HƯỚNG DẪN SỬ DỤNG SẢN PHẨM THI CÔNG}
% Nội dung phụ lục B

% ================= CHỈ MỤC =================
\clearpage
\phantomsection
\addcontentsline{toc}{chapter}{CHỈ MỤC}
\begingroup\raggedright
\printindex
\endgroup

% ================= TÀI LIỆU THAM KHẢO (kiểu IEEE) =================
% Chỉ liệt kê tài liệu thực sự được trích dẫn trong bài, không liệt kê tài liệu chỉ đọc.
\clearpage
\phantomsection
\addcontentsline{toc}{chapter}{TÀI LIỆU THAM KHẢO}
\bibliographystyle{IEEEtran}
\bibliography{References}

\end{document}
